% !TeX program = lualatex
\documentclass[11pt,a4paper]{ltjsarticle}
\usepackage[margin=25mm,headheight=17pt]{geometry}
\usepackage[haranoaji]{luatexja-preset}
\usepackage{amsmath,amssymb,amsthm,mathtools,mathrsfs,booktabs,array,enumitem,longtable}
\usepackage[unicode,hidelinks]{hyperref}
\usepackage{fancyhdr}
\pagestyle{fancy}\fancyhf{}\fancyfoot[C]{\thepage}
\fancyhead[L]{\small\nouppercase{\leftmark}}
\renewcommand{\sectionmark}[1]{\markboth{#1}{}}
\setlength{\emergencystretch}{3em}\allowdisplaybreaks[2]
\setlength{\parskip}{.15em}
\newtheoremstyle{jp}{.6em}{.6em}{\normalfont}{}{\bfseries}{.}{.5em}{}
\theoremstyle{jp}
\newtheorem{theorem}{定理}[section]
\newtheorem{proposition}[theorem]{命題}
\newtheorem{lemma}[theorem]{補題}
\newtheorem{corollary}[theorem]{系}
\newtheorem{definition}[theorem]{定義}
\newtheorem{remark}[theorem]{注意}
\newtheorem{example}[theorem]{例}
\newtheorem{question}[theorem]{問題}
\renewcommand{\proofname}{証明}
\newcommand{\kk}{\mathbb C}\newcommand{\PP}{\mathbb P}
\newcommand{\OO}{\mathcal O}\newcommand{\A}{\mathcal A}
\newcommand{\B}{B_{\mathrm{VK}}}\newcommand{\Blo}{B_{p,q}}
\newcommand{\HH}{\operatorname{HH}}\newcommand{\Hom}{\operatorname{Hom}}
\newcommand{\End}{\operatorname{End}}\newcommand{\Sym}{\operatorname{Sym}}
\newcommand{\Perf}{\operatorname{Perf}}\newcommand{\Der}{\operatorname{Der}}
\newcommand{\Inn}{\operatorname{Inn}}\newcommand{\gr}{\operatorname{gr}}
\newcommand{\rk}{\operatorname{rank}}\newcommand{\Spec}{\operatorname{Spec}}
\newcommand{\id}{\operatorname{id}}\newcommand{\Ext}{\operatorname{Ext}}
\newcommand{\Tor}{\operatorname{Tor}}\newcommand{\Ann}{\operatorname{Ann}}
\newcommand{\coker}{\operatorname{coker}}\newcommand{\Supp}{\operatorname{Supp}}
\newcommand{\pd}{\operatorname{pd}}\newcommand{\depth}{\operatorname{depth}}
\newcommand{\Frac}{\operatorname{Frac}}\newcommand{\Aut}{\operatorname{Aut}}
\newcommand{\Out}{\operatorname{Out}}\newcommand{\GL}{\operatorname{GL}}
\newcommand{\PGL}{\operatorname{PGL}}\newcommand{\Hilb}{\operatorname{Hilb}}
\newcommand{\dashsim}{\mathrel{\overset{\mathrm{bir}}{\sim}}}
\setcounter{tocdepth}{1}

\hypersetup{pdftitle={重み付きCalabi--Yau量子化の中心と表現}}
\title{重み付きCalabi--Yau量子化の中心と表現\\
\large 全中心、単純加群の合流と完備局所Brauer障害}
\author{}\date{2026年10月30日\quad 日本語公表用原稿XI・v40}
\begin{document}\maketitle
\begin{abstract}
重み $(1,1,1,m)$ の通常Calabi--Yau量子化について、中心と表現の幾何を調べる。
正共鳴の非PI部分では原始スペクトルを代数的なPoisson葉の空間と同定し、
中心と葉の閉包から全域の位相を記述する。
端点の一の冪根量子化では、混合項が $H=f(Q,u)+xG$ となる族の全中心を
二つの明示方程式による四次元正規Gorenstein完全交叉として求める。
$f=0$ ではAzumaya部分が滑らかな部分と一致し、全閉中心ファイバーを分類できる。
最深点では唯一の単純加群が一次元となり、根基の冪零指数は $2\ell-1$ であるが、
完備局所分数体上のBrauer周期と指数はともに $\ell$ を保つ。
従って、この完成代数は反射的加群の自己準同型環としてのNCCRへはならない。
$f\in Q\mathbb C[Q,u]$ でも、この最深ファイバーと完備局所指数は変わらない。
重み四・位数五の $f=0$ と $f=Qu^3$ は同じ関数斜体を持つ一方、
全中心の特異部分の次元が二と一に分かれ、アフィン代数としてdg導来Morita非同値となる。
一般の混合項に対しては、別の中心開部分を巡回ノルムによるアフィン改変として記述し、
標識付き境界から混合項の剰余類を回収する。
その巡回安定化群の位数が $g$ なら、一般境界の横断中心は $A_{g-1}$ 型であり、
複素閉点完成はMcKay型NCCRを持つ。
一方、一般境界点の可換解消の第 $i$ 例外曲線には位数 $g/\gcd(g,i)$ のBrauer留数が残る。
留数を消す最小係数被覆は、境界の単純加群を分離する巡回被覆と一致する。
\end{abstract}
\tableofcontents\clearpage
\section{序論}

中心の各閉点で量子化を理解できても、その局所表示を境界全体で貼り合わせられるとは限らない。
本稿では全中心と単純加群の合流を求め、さらに中心の解消・完備化に残るBrauer障害を計算する。
特に最深点では、中心ファイバーの単純加群が一次元一つになる一方、
完備局所分数体上のBrauer指数は量子化のPI次数と同じ最大値を保つ。

重み付き射影空間の非可換座標環には、射影圏の変形だけでは見えない
中心と表現の幾何がある。本稿では、非線形な局所冪零導分から作る
四次元CGG量子化族について、この幾何を正規元の局所化により具体的に求める。
特に重み $m=2j+1\ge3$ の $r=1$ 共鳴のCalabi--Yau代表では、
中心に $D_{m+4}$ 型曲面が現れ、その原点の上に非自明なprimitive spectrumが生じる。

まず正共鳴の非PI部分について、主要な結論を三点にまとめる。
第一に、正規元 $x$ を反転した環の中心は $A_{m+3}$ 型であり、
元の $D_{m+4}$ 型中心から一つの例外曲線 $\mathbb P^1$ を取り出す
crepantな部分解消の一つのchartになる。
第二に、中心の原点以外の全てのcentral reductionは
$\mathcal D(\mathbb G_m)$ に同型である。中心の分数体上でも根の付加を必要としない。
第三に、全てのprimitive idealを分類し、中心原点上では
\[
 \operatorname{Prim}_0(\mathcal S)
 =\{\mathfrak p_\lambda:\lambda\in\mathbb P^1(\mathbb C)\}
 \ \sqcup\
 \{\mathfrak n_{z_0,w_0}:(z_0,w_0)\in\mathbb A^2(\mathbb C)\}
\]
という表示を得る。
後者は全ての有限次元単純加群から来る文字表現の族である。
本稿が分類するのは単純加群の零化イデアルであり、
同じ零化イデアルを持つ無限次元単純加群まで同型類ごとに列挙するものではない。

原点上の閉包関係も決定できる。
$\mathfrak p_\lambda$ の一点閉包には、その点自身に加えて
\[
 \{\mathfrak n_{z_0,0}:z_0\in\mathbb C\}
\]
という共通の直線が含まれる。
一方、$\mathfrak p_\lambda$ の族の任意の無限部分集合は、原点上のprimitive spectrum全体に稠密である。
従って、$w_0\ne0$ の文字表現は一つの $\mathfrak p_\lambda$ の特殊化ではないが、
無限個の $\lambda$ を合わせると閉包に現れる。
この閉包公式によって、原点fiber内の任意の部分集合の位相を記述できる。

さらに全ての正共鳴で、量子素スペクトルとPoisson素スペクトルを同じ部分空間として同定する。
原始スペクトルは代数的な葉の空間であり、葉の飽和を取ってからZariski閉包を取る公式が、
非原点を動く中心指標の集積も含めて全域的な位相を記述する。
$r=1$ の原点の根基については、正確な冪零指数 $(m+7)/2$ も得る。
このスキーム論的な厚みは、原始スペクトルの位相とは別の情報である。

\subsection{方法と先行研究}

CGG積を使ったAS正則代数の構成と、正規元局所化によるスペクトルの研究には、
Lecoutre--Sierra\cite{LecoutreSierra2019}の先行研究がある。
彼らの線形な局所冪零導分による族と、本稿の非線形共鳴族は異なるが、
Ore拡大と可解Lie作用を使う方法を共有する。
本稿では局所化中心の指数半群、明示的な逆座標変換、
可換saturationと星積イデアルの一致を組み合わせる。

局所化後の単純加群を元の環へ制限しても、単純性は一般に保存されない。
この点は、可算次元複素整域のJacobson根基が零であることを使う
正規局所化の補題によって処理する。
境界の包含関係には、局所化した曲線の有限整拡大を明示してlying-overを適用する。
これにより、$x\ne0$ の計算だけから境界を推測せずに済む。

全primitive idealの表示から、Dixmier--Moeglin同値も導く。
これは本例のrational primeが全て明示したリストに入り、
正規局所化により局所閉になることによる。
一般のNullstellensatzの含意についてはBell--Wang--Yee\cite{BellWangYee2019}を参照する。

本稿の量子化族のAS正則性と全中心の生成性は、公表稿IX\cite{VKPoisson34}を入力とする。
射影圏の導来分類と変形環は、公表稿X\cite{VKCategories36}にまとめた。
ここでは中心の局所化、部分解消、central reduction、全正共鳴のスペクトル対応と原始イデアルの証明を
独立した一つの論証として記述する。

\subsection{範囲について}

この正共鳴の非PI部分で扱う代数は中心上有限ではなく、通常のAzumaya代数ではない。
中心がDu Val曲面であることからNCCRや非可換解消を自動的に得るとも主張しない。
部分解消の例外曲線と $\mathfrak p_\lambda$ のパラメータは共に $\mathbb P^1$ となり、
有限座標は局所中心の $c=qJ/x$ によって一致する。
しかし部分解消全体の上に代数の層を構成することは、ここでのprimitive idealの分類とは別の課題である。

\subsection{端点のPIファイバーとMcKay模型}

偶数端点の全変形を実現するAS正則CY族では、一般のパラメータは非PIである一方、
一の冪根では任意に高いPI次数を持つファイバーが現れる。
その中心の $x^\ell\ne0$ という開部分は、混合項を消去した一般化Weyl代数で計算できる。
位数 $\ell$ が奇数なら中心は滑らかで全体がAzumaya、
偶数 $2d$ なら $A_1$ 型の特異部分が正確に非Azumaya部分となる。
特異中心ファイバーは、二頂点と各向き二本の矢を持つ平方零代数の $d$ 次行列環である。

この表現の構造を、中心の有限エタール被覆上での滑らかな四次元環と
$\mu_2$ 作用へ持ち上げる。これにより、ファイバーだけでなく完備近傍全体が
McKay skew group模型で説明できる。
この完成局所表示が元の中心全体での自己準同型環表示を意味するかを調べると、
明確な障害がある。一般点の代数は
\[
 D=(\Xi,U)_t,\qquad
 x^\ell=\Xi,\quad u^\ell=U,\quad ux=txu
\]
という巡回斜体で、そのBrauer周期と指数はいずれも $\ell$ である。
従って、元の中心上で有限反射的加群の自己準同型環になることも、そのような環と
中心上でMorita同値になることもない。
完成局所的なNCCRと大域的な自己準同型環の相違を、この周期が検出する。

また特異部分 $\Sigma\simeq\mathbb G_{m,Q}\times\mathbb A^1_J$ 上の全代数は、
$L=\mathbb C[Q^{\pm1},J,\kappa]/(\kappa^2+tQ/\beta)$ による
\[
 A_\Sigma\simeq M_d\bigl(L\ltimes(L_\iota\oplus L_\iota)\bigr),
 \qquad\iota(\kappa)=-\kappa
\]
である。二つの単純加群はこの連結二重被覆上で分離し、$Q$ の一周によって交換される。
式は特異部分への制限全体について成り立ち、閉点ごとの模型だけに留まらない。

\subsection{混合項を回収する境界と巡回被覆}

関数斜体と $x^\ell\ne0$ の中心からは消えていた $H$ を検出するには、
別の中心開部分 $U=u^\ell\ne0$ が有効である。
ここで $H=h(Q,u)+xG$ と分けて $W=w+G$ と置くと、
\[
 xW=J-h(Q,u),\qquad Wx=J-h(Q,tu)
\]
という一般化Weyl表示となる。その全中心は
\[
 \mathbb C[Q,J,U^{\pm1},\Xi,\Omega]/(\Xi\Omega-P_h),
 \qquad P_h=\prod_{i=0}^{\ell-1}(J-h(Q,t^iu))
\]
であり、$\Omega=W^\ell$ である。
従って $H$ は中心の境界への延長を変える。
$Q,U,J$ を固定して $P_h$ を知れば、$h$ は $u\mapsto t^ku$ まで決まる。
重み四・位数五では、$H=u^5$ と $H=u^4v$ が同じ関数斜体と同じ
$x^\ell\ne0$ の開部分を持ちながら、$U\ne0$ の中心の特異部分の次元は異なる。

境界ファイバーを決めるのは $J=h(Q,t^i\lambda)$ が成立する巡回軌道上の位置である。
成立する辺で巡回箙を切り、残る鎖の長さを $d_1,\ldots,d_r$ とすると、
単純加群の次元は正確に $d_1,\ldots,d_r$ であり、根基の冪零指数は $r$ となる。
さらに $h$ の巡回安定化群が位数 $g$ なら、衝突を除いた正規化境界上の
全代数と、連結巡回 $g$ 重被覆上の単純表現族を明示できる。
この部分でもAzumaya部分は中心の滑らかな部分と一致する。

一般化Weyl代数の重み加群と巡回軌道の理論はDrozd--Guzner--Ovsienko
\cite{DGOBoundary39}に帰する。
一の冪根での量子代数の中心・Azumaya部分の具体的研究として
Cooney--Ganev--Jordan\cite{CGJBoundary39}も参照する。
本稿では、この枠組みを混合Ore族の境界へ適用し、全ファイバーの行列表示、
標識付きノルム復元、および元の中心上への降下を明示する。
巡回代数のノルム分裂条件\cite{MinacSchultzSwallowDescent39}から得るBrauer障害と、
境界の有限被覆が記録する情報を、次の形式近傍の計算によって結び付ける。

\subsection{全中心による区別と最深点に残る指数}

混合項が $H=f(Q,u)+xG$ と書ける場合には、$W=w+G$ と直し、
$\Omega=W^\ell$ を加えることで全中心を得る。
巡回ノルムを $N_f(Q,U,J)=\prod_i(J-f(Q,t^iu))$ とすると、
\[
 Z(S_H)=\frac{\mathbb C[Q,J,\Xi,\Omega,U,V]}
 {(UV-P_\ell(Q,\Xi),\ \Xi\Omega-N_f(Q,U,J))}
\]
である。正規な二つの開部分だけでなく、残る余次元二の部分も処理することで、
この表示が局所化前の全中心であることを証明する。

重み四、$\ell=5$、$\beta=1$ で $f=0$ と $f=Qu^3$ を比べると、
第二の方程式はそれぞれ $\Xi\Omega=J^5$ と
$\Xi\Omega=J^5-Q^5U^3$ となる。
前者の特異部分は二次元成分を持つが、後者の被約特異部分は三本の座標軸である。
従って標識を忘れた全中心も非同型であり、二つのアフィン代数は
Morita同値でも $\mathbb C$ 線形dg導来Morita同値でもない。
両者の関数斜体と $\Xi\ne0$ の中心開部分は同じである。
射影圏の同値を排除するには、さらに次数付きの不変量が必要となる。

混合項が $H=H_c+xG$、$H_c\in Z(B)$ と書ける場合には、$w\mapsto w+G$ により
$H=0$ の部分族へ移る。その全中心は
\[
 Z_0=\frac{\mathbb C[Q,J,\Xi,\Omega,U,V]}
 {(\Xi\Omega-J^\ell,\ UV-P_\ell(Q,\Xi))}
\]
であり、$\ell$ の偶奇により
\[
 P_\ell=-Q^\ell-\beta^\ell\Xi^2\quad(\ell\text{ 奇数}),\qquad
 P_{2d}=(Q^d+(-1)^d\beta^d\Xi)^2
\]
となる。これは四次元の正規Gorenstein完全交叉である。
全ての閉中心ファイバーを分類し、Azumaya部分が中心の滑らかな部分と一致すると示す。
最深点では一つの一次元単純加群しか残らないが、ファイバーの全次元は
奇数時 $(7\ell^2-6\ell+1)/2$、偶数時 $(7\ell^2-6\ell+2)/2$ へ増大する。
根基の冪零指数は正確に $2\ell-1$ である。

この原点で中心を完備化しても、一般点の巡回代数の周期と指数は $\ell$ のままである。
証明では $\operatorname{wt}(Q,J,\Xi,U,V)=1$、$\operatorname{wt}(\Omega)=\ell-1$ の
付値を作り、剰余体の独立変数 $U/\Xi$ がノルムであることを妨げる。
従ってこの完成代数は、中心上の反射的加群の自己準同型環とも、その中心線形Morita型ともならない。
$\Xi\ne0$ の完成局所McKay模型を、最深点にまで外挿することはできない。
$f$ が $Q$ の倍数である場合にも、最深点のファイバーと完備局所指数は同じである。
従って上記の全中心が非同型な二例は、これらの最深点の数値も共有している。

\subsection{例外因子の留数としての巡回モノドロミー}

一般の混合項についても、$U\ne0$ の一般境界では全形式近傍を記述できる。
$h$ の安定化群位数を $g$、$n=\ell/g$、$z=u^g$ とすると、
正規化境界の関数体 $F=\mathbb C(Q,z)$ 上の横断中心は
\[
 R_F=F[[\xi,\upsilon,s]]/(\xi\upsilon-s^g)
\]
となる。完成代数は巡回拡大 $L=F[\kappa]/(\kappa^g-z)$ を係数とする角の $n$ 次行列環であり、
$L$ への拡大後には $L[[a,b]]\#\mu_g$ という標準McKay模型になる。
特に一般境界の複素閉点完成はNCCRを持つ。

しかし一般境界点の最小可換解消では、第 $i$ 例外曲線に
\[
 \partial_{E_i}[D]=i\chi,\qquad
 \operatorname{ord}\partial_{E_i}[D]=g/\gcd(g,i)
\]
というBrauer留数が残る。ここで $\chi$ は $L/F$ に対応する巡回指標である。
この式は標準的な巡回代数の留数公式\cite[Proposition~5.8]{HarariBrauer40}の本族への適用である。
全ての留数を消す最小有限係数被覆は $L/F$ であり、
単純加群の巡回モノドロミーを解く被覆と一致する。
これにより、閉点完成のNCCR、一般境界での分岐、最深点に残る指数を一つの計算体系で比較できる。

\section{量子化族と多項式不変量}
\label{sec:primitive-setup-v36}

基礎体を $\mathbb C$ とする。整数 $j,r\ge0$、$m=2j+r\ge3$ に対して
\[
 R_m=\mathbb C[x,y,z,w],\quad \deg(x,y,z,w)=(1,1,1,m),\quad
 q=xz-y^2,\quad d=m+r+1
\]
と置く。次の二つの導分を使う。
\begin{align*}
 N&=x\partial_y+2y\partial_z-q^jz^r\partial_w,\\
 X_\alpha&=\alpha x\partial_x+(\alpha+1)y\partial_y+(\alpha+2)z\partial_z
                  +(m\alpha+d)w\partial_w.
\end{align*}
$N$ は局所冪零、$[N,X_\alpha]=N$ である。
同じベクトル空間に有限積
\[
 a\star b=\sum_{n\ge0}s^nN^n(a)\binom{X_\alpha}{n}(b)
\]
を入れた代数を $\mathcal S_{j,r}(\alpha,s)$ または
$\mathcal S^{(q)}_{j,r}(\alpha,s)$ と書く。
積の結合性はCGG恒等式による\cite[Proposition~2.2]{LecoutreSierra2019}。
本稿では量子化係数 $s\ne0$ を固定する。

\begin{proposition}[基本的な入力]
\label{prop:primitive-ring-input-v36}
$\mathcal S$ はnoetherianなAS正則四次元整域であり、
\[
 H_{\mathcal S}(t)=\frac1{(1-t)^3(1-t^m)}.
\]
$T=1+sN$ とその全スカラー冪は自己同型であり、Nakayama自己同型は
$T^{(m+3)\alpha+m+r+3}$ である。
従って通常のCalabi--Yau代表は
$\alpha=-(m+r+3)/(m+3)$ で与えられる。
\end{proposition}

これらの性質と以下の不変環表示は、公表稿IX\cite{VKPoisson34}の反復Ore表示と
多項式不変量の計算による。量子化の発散だけからCalabi--Yau性を推論しない。
\[
 c_r=(-1)^r\frac{(2r)!!}{(2r+1)!!},\qquad
 J=x^{r+1}w+q^j\sum_{k=0}^r\binom rk\frac{q^{r-k}y^{2k+1}}{2k+1},\qquad
 \Delta=\frac{J^2+c_r^2q^d}{x^{r+1}}
\]
と置くと、$\Delta$ は多項式であり、
\[
 \ker N=\mathbb C[x,q,J,\Delta]
 \simeq\mathbb C[X,Q,U,V]/(X^{r+1}V-U^2-c_r^2Q^d).
\]
$N$ 不変元同士の星積は通常の可換積と等しい。
通常CY値での全中心は
\[
 Z(\mathcal S)=
 \begin{cases}
 \mathbb C&r\ge2,\\
 \mathbb C[q,J]&r=0,\\
 \mathbb C[\mathsf S,\mathsf P,\mathsf R]/
 (\mathsf R^2-\mathsf S\mathsf P^2-\frac{16}{81}\mathsf S^{m+3})&r=1,
 \end{cases}
\]
ただし最後の場合は
\[
 \mathsf S=q^2\Delta,\quad
 \mathsf P=2xJ\Delta^{(m+3)/2},\quad
 \mathsf R=q(x^2\Delta+J^2)\Delta^{(m+3)/2}
\]
である。この中心の全体性には、正規元 $x$ による飽和と零重量単項式の分類を使う。

以後、$\mathcal S[x^{-1}]$ の中心、射影的次数零chartの中心、
および $Z(\mathcal S)$ は区別して記す。
また、primitive idealとは単純左加群の零化イデアルを指し、
$\operatorname{Prim}\mathcal S$ には両側イデアルの包含で定めるJacobson位相を入れる。
本稿の全primitive idealは完全素となるので、各商環の分数斜体を通常のOre斜体として扱える。

\section{正規元の局所化とDu Val中心の特殊化}
\label{sec:central-reductions-v35}

この節では、前節で構成した量子代数の中心を、単なる関数体の情報より詳しく調べる。
正規元 $x$ を反転すると中心には $A$ 型特異点が現れ、
$r=1$ では元の $D$ 型中心との間に明示的なblowup chartが得られる。
さらに $D$ 型中心の原点以外の全central reductionを決定する。
以下、基礎体は $\mathbb C$、量子化係数は $s\ne0$ とする。

\subsection{全共鳴の局所化とその中心}

$j,r\ge0$、$m=2j+r\ge3$、$d=m+r+1$ とし、
$\mathcal S=\mathcal S_{j,r}^{(q)}(\alpha,s)$ と書く。
$q,J$ は前節の $N$ 不変多項式である。

\begin{proposition}[Euler型の局所表示]
\label{prop:euler-localization-v35}
$x$ の冪は両側Ore集合をなし、全ての $\alpha$ について
\[
 \mathcal S[x^{-1}]
 \simeq \mathbb C[x^{\pm1},q,J][t;sX_\alpha],
 \qquad t=x^{-1}\star y,
\]
となる。右辺の導分は
\[
 X_\alpha(x)=\alpha x,\qquad
 X_\alpha(q)=2(\alpha+1)q,\qquad
 X_\alpha(J)=d(\alpha+1)J
\]
で与えられる。
\end{proposition}

\begin{proof}
$T=1+sN$ とすると
$a\star x=T^\alpha(a)x$、$x\star a=xa$ である。
$T^\alpha$ は自己同型で $x$ を固定するので、$x$ は正則正規元である。
$x$ の逆を加えた多項式上にも $N$ は局所冪零に延び、同じ有限星積を使える。
$N$ 不変元同士の星積は通常積であり、
\[
 [y,f]_\star=sxX_\alpha(f)\qquad(N(f)=0)
\]
だから $[t,f]=sX_\alpha(f)$ を得る。
通常の可換座標変換では
\[
 \mathbb C[x^{\pm1},y,z,w]
 =\mathbb C[x^{\pm1},q,J,t],\qquad t=y/x.
\]
実際 $z=(q+y^2)/x$ であり、$J$ は $w$ に関して可逆係数 $x^{r+1}$ を持つ。
また $N(t)=1$、$X_\alpha(t)=t$ なので、星積による $t$ の冪は通常冪に対して
最高係数一の三角変換になる。
従って右辺のOre環の基底 $1,t,t^2,\ldots$ は左辺でも基底となり、同型が従う。
\end{proof}

\begin{theorem}[局所化した中心の $A$ 型表示]
\label{thm:localized-center-v35}
通常のCalabi--Yau値
\[
 \alpha=-\frac{m+r+3}{m+3}
\]
を取る。$A=m+r+3=d+2$、$g=\gcd(A,r)$、$n=A/g$ と置くと
\[
 \boxed{
 Z(\mathcal S[x^{-1}])
 =\frac{\mathbb C[a,b,c]}{(ab-c^n)},\quad
 a=\frac{q^n}{x^{2r/g}},\quad
 b=\frac{J^n}{x^{rd/g}},\quad
 c=\frac{qJ}{x^r}.
 }
\]
$r=0$ では $g=A,n=1$ と解釈し、これは $\mathbb C[q,J]$ である。
$r>0$ の場合は $A_{n-1}$ 型の曲面となる。
\end{theorem}

\begin{proof}
係数環 $R=\mathbb C[x^{\pm1},q,J]$ 上の導分は
\[
 sX_\alpha=-\frac{s}{m+3}
       (Ax\partial_x+2r q\partial_q+rdJ\partial_J).
\]
$t$ 次数が正の元の $x$ との交換子は最高次数項を一つ下げた非零項を持つ。
従ってOre環の中心は $R$ の導分核に等しい。
$r>0$ とする。零重量単項式 $x^e q^uJ^v$ は
\[
 Ae+2ru+rdv=0,\qquad e\in\mathbb Z,\quad u,v\ge0
\]
を満たす。$A=d+2$ は奇数で、$\gcd(A/g,r/g)=1$ だから、
この条件は $u\equiv v\pmod n$ と同値である。
従って $c^{\min(u,v)}$ を取り出すと、残りは $a$ または $b$ の冪になる。
関係は $ab=c^n$ であり、両方の $a,b$ を含まない単項式の像は
互いに異なるLaurent単項式なので、これ以外の関係はない。
$r=0$ では重量条件は単に $e=0$ となる。
\end{proof}

通常の重み付き次数では
\[
 \deg(a,b,c)=\left(\frac{2(m+3)}g,
                    \frac{d(m+3)}g,m+3\right).
\]
ここで扱ったのは全次数付き環の局所化である。
射影的な $x\ne0$ chartの次数零環は、別に
\[
 (\mathcal S[x^{-1}])_0
 =\mathbb C[u,v][t;s(2u\partial_u+dv\partial_v)],
 \qquad u=q/x^2,\quad v=J/x^d
\]
となり、その中心は $\mathbb C$ である。
また $x$ は中心元ではないため、前定理は中心の通常の可換局所化を主張していない。
例えば $r\ge2$ では $Z(\mathcal S)=\mathbb C$ であるのに、
$Z(\mathcal S[x^{-1}])$ は二次元環である。

\subsection{$D$ 型中心から $A$ 型中心へのblowup chart}

以下では $r=1$、$m=2j+1\ge3$ とする。
\[
 d=m+2,\quad A=d+2=m+4,\quad
 L=(d+1)/2=(m+3)/2,\quad \beta=4/9
\]
と置く。元の中心の生成元を $\mathsf S,\mathsf P,\mathsf R$ と書けば
\[
 Z=Z(\mathcal S)
 =\frac{\mathbb C[\mathsf S,\mathsf P,\mathsf R]}
 {\bigl(\mathsf R^2-\mathsf S\mathsf P^2-\beta^2\mathsf S^{2L}\bigr)}
\]
であり、$D_A$ 型である。

\begin{proposition}[中心間の明示的な双有理射]
\label{prop:blowup-chart-v35}
$B=Z(\mathcal S[x^{-1}])=\mathbb C[a,b,c]/(ab-c^A)$ とすると、
$Z\subset B$ は次で与えられる。
\begin{align*}
 \mathsf S&=c^2+\beta a,\\
 \mathsf P&=2\left(b+
   \sum_{k=0}^{L-1}\binom Lk\beta^{L-k}
                   a^{L-1-k}c^{2k+1}\right),\\
 \mathsf R&=c\mathsf P+\beta\mathsf S^L.
\end{align*}
さらに共通の関数体の中で
\[
 \boxed{
 B=Z\left[\frac{\mathsf R-\beta\mathsf S^L}{\mathsf P}\right].
 }
\]
従って $\operatorname{Spec}B$ は、イデアル
$(\mathsf P,\mathsf R-\beta\mathsf S^L)$ のblowupの $\mathsf P$-chartである。
原点のスキーム論的逆像は $\mathbb A^1_c$ である。
\end{proposition}

\begin{proof}
$\Delta=(J^2+\beta q^d)/x^2$ と前節の中心生成元を用いると
\[
 \mathsf S=c^2+\beta a,\qquad
 a\mathsf P=2c\mathsf S^L,\qquad
 a\mathsf R=(\mathsf S+c^2)\mathsf S^L
\]
が得られる。最初の二式を展開し $ab=c^A$ を使うと、表示した多項式公式になる。
三式目から $\mathsf R-c\mathsf P=\beta\mathsf S^L$ が従う。
従って $c=(\mathsf R-\beta\mathsf S^L)/\mathsf P$、
$a=(\mathsf S-c^2)/\beta$ であり、$b$ も $\mathsf P$ の式から復元できる。
これで部分環としての等式を得る。

blowup chartの定義は $Z[I/\mathsf P]\subset Z[\mathsf P^{-1}]$ なので、
上の部分環等式は余分な $\mathsf P$-torsionを含まないchartを直接同定している。
実際、$\mathsf R=c\mathsf P+\beta\mathsf S^L$ を代入してから
$\mathsf P$ 因子を除くと
\[
 (\mathsf S-c^2)\mathsf P=2\beta c\mathsf S^L
\]
となる。$a=(\mathsf S-c^2)/\beta$ と置き、表示した $b$ へ変数変換すれば
この式は正確に $2\beta(ab-c^A)=0$ となる。

原点上では $\mathsf S=0$ から $a=-c^2/\beta$、
$\mathsf P=0$ から $b=-\beta c^d$ を得る。
残る関係は全て恒等式となるので、そのファイバーは $\mathbb C[c]$ である。
\end{proof}

このchartは $D_A$ の全体の解消ではない。
曲線 $\mathsf P=0,\mathsf R=-\beta\mathsf S^L$ の原点以外には逆像を持たない。
一方、$\mathsf P\ne0$ 上では上の $c$ の式で逆射を持ち、
$\mathsf R+\beta\mathsf S^L\ne0$ 上では
\[
 c=\frac{\mathsf S\mathsf P}{\mathsf R+\beta\mathsf S^L},\qquad
 a=\frac{2\mathsf S^{L+1}}{\mathsf R+\beta\mathsf S^L}
\]
で逆射を持つ。

\begin{corollary}[一曲線を取り出すcrepantな部分解消]
\label{cor:partial-resolution-v35}
上のイデアルのblowup全体を $Y\to\operatorname{Spec}Z$ とする。
これは原点以外で同型であり、原点の逆像は一つの $\mathbb P^1$ である。
$Y$ は正規Gorenstein曲面で、残る唯一の特異点は $A_{A-1}$ 型である。
またこの射はcrepantである。
\end{corollary}

\begin{proof}
残るchartで $U=\mathsf R-\beta\mathsf S^L$、$v=\mathsf P/U$ と置く。
$U$-torsionを除いた方程式は
\[
 U(1-\mathsf Sv^2)+2\beta\mathsf S^L=0
\]
となる。この環では $1-\mathsf Sv^2$ は単元である。
実際 $z=\mathsf Sv^2$ と書き、
$1=(1-z)(1+z+\cdots+z^{L-1})+z^L$ に上の式を代入すれば単元の逆を得る。
従ってこのchartは
$\mathbb C[\mathsf S,v,(1-\mathsf Sv^2)^{-1}]$ であり滑らかである。
原点上の $\mathbb A^1_v$ と前chartの $\mathbb A^1_c$ は $c=1/v$ で貼り合わさる。
原点以外では中心イデアルは可逆である。実際
$\mathsf R+\beta\mathsf S^L$ が可逆な点では
$U=\mathsf S\mathsf P^2/(\mathsf R+\beta\mathsf S^L)$ なのでイデアルは $(\mathsf P)$、
それ以外の非原点でイデアルが消えることはない。
よってblowupは原点以外で同型である。

前chartの $A_{A-1}$ 表示と後chartの滑らかさから正規Gorenstein性も従う。
正則二形式を比較すると、有理形式として
\[
 \frac{d\mathsf S\wedge d\mathsf P}{2\mathsf R}
       =\frac{da\wedge dc}{a}
\]
である。これは $\mathsf S=c^2+\beta a$、
$\mathsf P=2c\mathsf S^L/a$ から直接計算できる。
右辺は $ab=c^A$ 上の双対化加群の生成元であり、例外曲線の一般点で零も極も持たない。
他に例外因子はないので相対標準因子は零となる。
\end{proof}

\subsection{原点以外の全central reduction}

$\mathcal D(\mathbb G_m)$ を
\[
 \mathbb C[\eta^{\pm1}][E;\eta\partial_\eta],\qquad
 [E,\eta]=\eta
\]
と書く。これはLaurent多項式環上の通常の微分作用素環である。

\begin{theorem}[滑らかな中心点における特殊化]
\label{thm:central-reductions-v35}
$r=1$、$m\ge3$ を奇数とし、通常のCY値と $s\ne0$ を取る。
$\chi\in\operatorname{Spec}Z$ の閉点が原点以外ならば
\[
 \boxed{\mathcal S/\mathfrak m_\chi\mathcal S
          \simeq\mathcal D(\mathbb G_m).}
\]
特にこれらのcentral reductionは単純整域である。
また $F=\operatorname{Frac}Z$ に対しても
\[
 \mathcal S\otimes_ZF
 \simeq F[\eta^{\pm1}][E;\eta\partial_\eta]
\]
が成立する。
\end{theorem}

\begin{proof}
$\chi=(S_0,P_0,R_0)$ と書く。$N$ 不変元同士の積は可換なので、
以下の恒等式は量子代数の中でもそのまま成り立つ。
\begin{align*}
 \mathsf S&=q^2\Delta,\qquad
 \mathsf P=2xJ\Delta^L,\\
 \mathsf R+\beta\mathsf S^L&=2q x^2\Delta^{L+1},\qquad
 \mathsf R-\beta\mathsf S^L=2q J^2\Delta^L,\\
 x^2\Delta&=J^2+\beta q^d.
\end{align*}
まず $P_0\ne0$ なら、可換する因子の積 $2xJ\Delta^L$ が可逆なので $x$ は可逆である。
$S_0\ne0$ かつ $R_0+\beta S_0^L\ne0$ でも、$q,\Delta$ と $x^2$ が可逆となり同じ結論を得る。

これらの場合は命題\ref{prop:euler-localization-v35}を適用できる。
中心 $B$ の $\chi$ 上のファイバーは一つの被約点 $(a_0,b_0,c_0)$ である。
$P_0\ne0$ なら $c_0=(R_0-\beta S_0^L)/P_0$ から順に復元できる。
$P_0=0$ の場合には $S_0\ne0$ であり、
$a\mathsf P=2c\mathsf S^L$ より $c_0=0$、$a_0=S_0/\beta$、$b_0=0$ を得る。
この点は $B$ の原点ではないので、$a_0,b_0$ の少なくとも一つが非零である。

$a_0\ne0$ とすると、係数環の特殊化は
\[
 \frac{\mathbb C[x^{\pm1},q^{\pm1},J]}
      {(q^A-a_0x^2,\ qJ-c_0x)}
 \simeq\mathbb C[\eta^{\pm1}],\qquad
 \eta=q^{(A+1)/2}/x.
\]
逆の式は
\[
 q=\eta^2/a_0,\qquad
 x=\eta^A/a_0^{(A+1)/2},\qquad
 J=c_0\eta^d/a_0^{(A-1)/2}
\]
である。$b_0\ne0$ の場合は $h=(d-1)/2$、
$\eta=J^{h+1}/x^h$ を使えば
\[
 x=\eta^A/b_0^{h+1},\qquad J=\eta^d/b_0^h,
 \qquad q=c_0x/J
\]
となる。いずれも $[t,\eta]=\sigma\eta$、
$\sigma=-s/(m+3)\ne0$ なので $E=t/\sigma$ と置けば所望の環になる。

残る非原点は
\[
 S_0\ne0,\qquad P_0=0,\qquad R_0=-\beta S_0^L
\]
を満たす。ここでは $q,\Delta$ が可逆であり、上の恒等式から $x^2=0$、
$J^2=-\beta q^d$ を得る。従って $J$ は可逆で、$\mathsf P=0$ より $x=0$ である。
$q=-y^2$ だから $y$ は可逆になる。
$x=0$ の商で
\[
 H=w+\frac{(-1)^j}{4}y^{m-2}z^2,\qquad
 \tau=(2y)^{-1}\star z
\]
と置けば $N(H)=0$、$N(\tau)=1$ であり、同じ三角基底の議論から
\[
 (\mathcal S/(x))[y^{-1}]
 =\mathbb C[y^{\pm1},H][\tau;sX_\alpha]
\]
となる。この表示では
\[
 X_\alpha(y)=-\frac{y}{m+3},\qquad
 X_\alpha(H)=\frac{m+6}{m+3}H,\qquad
 \mathsf S=-\frac43(-1)^j y^{m+6}H.
\]
$\mathsf S=S_0$ により $H$ を消去できる。他の二つの中心式は自動的に成立する。
従って $\eta=y$、$E=\tau/\sigma$ とすれば再び $\mathcal D(\mathbb G_m)$ を得る。

$F$ 上の一般ファイバーは最初の $a_0\ne0$ の計算と同じである。
表示した式に定数の冪根は不要なので、代数閉包への拡大も必要ない。

最後に、$\mathcal D(\mathbb G_m)$ の非零両側イデアルの元で $E$ 次数最小のものを選び、
$\eta$ との交換子を取ると、その最小次数は零となる。
従ってイデアルは $\mathbb C[\eta^{\pm1}]$ と非自明に交わる。
その交わりは $\eta\partial_\eta$ で安定であり、単項式の個数最小の元に
$\eta\partial_\eta-\lambda$ を作用させると、単項式を含むことが分かる。
Laurent単項式は可逆なのでイデアルは全体である。整域性はOre表示から従う。
\end{proof}

微分Ore環の単純性の一般理論は\cite{ORS2013}、
Ore環とPoisson素スペクトルの比較は\cite{JordanOre2014}を参照されたい。
ここでは本族に必要な単純性を直接証明した。
上の定理は原点のcentral reductionを単純とする主張でも、中心上の有限性を主張するものでもない。

\subsection{原点の局所化ファイバーとそのprimitive ideal}

元の中心原点の極大イデアルを $\mathfrak m_0=(\mathsf S,\mathsf P,\mathsf R)$ とする。
この項では原代数の商全体ではなく、$x$ を反転した商を扱う。
\[
 h_0=J^2+\beta q^d,\qquad
 \delta=\sigma(Ax\partial_x+2q\partial_q+dJ\partial_J),
 \quad \sigma=-s/(m+3)
\]
と置く。

\begin{proposition}[原点に現れる局所化後の厚み]
\label{prop:localized-origin-v35}
$\mathcal B_0=(\mathcal S/\mathfrak m_0\mathcal S)[x^{-1}]$ には表示
\[
 \boxed{
 \mathcal B_0=
 \left(\frac{\mathbb C[x^{\pm1},q,J]}
                    {(q^2h_0,Jh_0^L)}\right)[t;\delta]
 }
\]
があり、その中心は $\mathbb C[c]$、$c=qJ/x$ である。
prime radicalは $(h_0)$、その冪零指数は正確に $L+1$ である。
その商は
\[
 \left(\frac{\mathbb C[x^{\pm1},q,J]}{(J^2+\beta q^d)}\right)[t;\delta]
\]
という整域になる。

$c=\lambda\ne0$ の中心ファイバーは $\mathcal D(\mathbb G_m)$ である。
一方、$c=0$ のファイバーは
\[
 \boxed{
 \mathcal B_0/(c)=
 \left(\frac{\mathbb C[x^{\pm1},q,J]}{(qJ,q^A,J^A)}\right)[t;\delta].
 }
\]
そのprime radicalは $(q,J)$、冪零指数は正確に $A$ であり、
商は $\mathcal D(\mathbb G_m)$ である。
このファイバーは部分環
$\mathbb C[x^{\pm1}][t;\sigma Ax\partial_x]\simeq\mathcal D(\mathbb G_m)$
上の左加群・右加群として、いずれも階数 $2A-1$ の自由加群である。
\end{proposition}

\begin{proof}
$x$ を反転すると $\Delta=x^{-2}h_0$ であり、
$\mathsf S=0,\mathsf P=0$ はそれぞれ $q^2h_0=0,Jh_0^L=0$ と同値になる。
$\mathsf R=c\mathsf P+\beta\mathsf S^L$ なので第三式は従う。
$x$ との交換子により中心元は係数環に入る。
$\delta$ は対角的であり、中心の零重量三式で割っているので、
商の零重量部分は $B/(\mathsf S,\mathsf P,\mathsf R)=\mathbb C[c]$ である。

恒等式
\[
 h_0^{L+1}=J(Jh_0^L)+\beta q^{d-2}h_0^{L-1}(q^2h_0)
\]
より $h_0^{L+1}=0$ である。
$q=0$ の商では $h_0^L$ は $J^{2L}$ となり、
関係は $J^{2L+1}=0$ だけなので $h_0^L\ne0$ である。
$d$ が奇数なので $J^2+\beta q^d$ は既約であり、
$h_0$ で割った係数環とそのOre拡大は整域となる。
従って $(h_0)$ はprime radicalそのものである。

$c=\lambda\ne0$ とすると $qJ=\lambda x$ から $q,J$ が可逆になり、$h_0=0$ を得る。
$h=(d-1)/2$、$\eta=J/q^h$ と置けば
\[
 q=-\eta^2/\beta,\qquad
 J=(-1/\beta)^h\eta^d,\qquad
 x=\lambda^{-1}(-1/\beta)^{h+1}\eta^A
\]
となる。また $[t,\eta]=\sigma\eta$ なので所望の微分作用素環である。

$c=0$ では $qJ=0$ であり、二つの中心式は正確に $q^A=0,J^A=0$ に帰着する。
従って係数環には
\[
 1,q,\ldots,q^{A-1},J,\ldots,J^{A-1}
\]
という $\mathbb C[x^{\pm1}]$ 基底がある。
この基底の各元は $\delta$ の固有元なので、Ore表示から左右の自由性が従う。
$(q,J)^A=0$、$(q,J)^{A-1}\ne0$ であり、その商は上記の単純な微分作用素環である。
従ってprime radicalの主張も従う。
\end{proof}

この結果から、$\mathcal B_0$ のprimitive idealは $\lambda\in\mathbb C$ によって
一意にパラメータ化される。元の局所化Ore環での逆像は
\[
 \lambda\ne0:\quad(h_0,qJ-\lambda x),\qquad
 \lambda=0:\quad(q,J)
\]
である。ここでは中心 $c$ が単純加群上でスカラーとなることを使う。
その理由は通常のDixmierの議論である：単純加群は可算次元の環の巡回商なので可算次元であり、
$c$ が超越的なら $(c-\lambda)^{-1}v$ が非可算個の一次独立なベクトルを与えてしまう。
代数的なら基礎体が代数閉体なのでスカラーとなる。
各スカラーに対するprimitive idealの一意性は前命題の単純商から従う。
この分類は、原代数の原点ファイバーの $x=0$ 側のprimitive idealを含まない。

\subsection{有限次元単純表現と中心原点}

\begin{proposition}[有限次元単純表現]
\label{prop:finite-simples-v35}
$r=1$、$m=2j+1\ge3$ の通常CY量子化では、全ての有限次元単純表現は一次元であり、
\[
 x=y=0,\qquad z=z_0,\quad w=w_0,
 \qquad (z_0,w_0)\in\mathbb A^2
\]
で与えられる。それらの中心指標は全て原点である。
\end{proposition}

\begin{proof}
有限次元単純加群上で正規元 $x$ は零または可逆に作用する。
可逆なら $t=x^{-1}y$ を用いた関係
$x^{-1}tx=t+s\alpha$ のtraceが矛盾するので、$x=0$ である。
$\alpha\ne0$ はCY値の式から分かる。
$\mathcal S/(x)$ では $y$ は正規元であり、可逆なら
$\tau=(2y)^{-1}z$ と $[\tau,y]=s(\alpha+1)y$ に同じtrace議論を適用できる。
$\alpha+1=-1/(m+3)\ne0$ なので $y=0$ となる。
$j\ge1$ だから、この二つを零にした商では $N$ が零となり、
商代数は通常の $\mathbb C[z,w]$ である。主張する全表現が得られる。
中心式では $q,J,\Delta$ がいずれも零になるので、中心指標は原点である。
\end{proof}

原点以外の各中心指標に対する両側イデアル $\mathfrak m_\chi\mathcal S$ は、
定理\ref{thm:central-reductions-v35}によりその指標を持つ唯一のprimitive idealである。
この一意性は同じ指標を持つ単純加群の一意性を意味しない。
原点上の全primitive idealとその閉包位相は次節で決定する。
原点のcentral reduction自体の全構造については、根基の正確な全域的冪零指数などが残る。


\begin{lemma}[原点のもう一つの正規局所化]
\label{lem:origin-y-chart-v36}
\[
 \left(\mathcal S/(\mathfrak m_0\mathcal S+(x))\right)[y^{-1}]
 \simeq\mathcal D(\mathbb G_m).
\]
\end{lemma}
\begin{proof}
定理\ref{thm:central-reductions-v35}の証明中の表示
$(\mathcal S/(x))[y^{-1}]=\mathbb C[y^{\pm1},H][\tau;sX_\alpha]$ を用いる。
この環では
\[
 \mathsf S=-\frac43(-1)^j y^{m+6}H,\qquad
 \mathsf P=0,\qquad \mathsf R=-\frac49\mathsf S^{(m+3)/2}.
\]
$y$ が可逆なので中心原点への特殊化は正確に $H=0$ を課す。
$\eta=y$、$E=-(m+3)\tau/s$ と置けば主張を得る。
ここで $y$ を反転するのは $x=0$ の商であり、元代数の $y$ の冪にOre条件は要求しない。
\end{proof}

\section{全primitive idealの分類と中心原点上の閉包}
\label{sec:primitive-spectrum-v36}

本節では $r=1$、$m=2j+1\ge3$、$s\ne0$、
$\alpha=-(m+4)/(m+3)$ とする。
$\mathcal S$ は前節のCGG量子代数であり、積を $\star$ と書く。
\[
 d=m+2=2j+3,\quad A=m+4=d+2,\quad L=(m+3)/2,
 \quad\varepsilon=(-1)^j,\quad\beta=4/9
\]
と置く。通常の可換多項式としての式
\begin{align*}
 q&=xz-y^2,\\
 J&=x^2w+q^j\left(xyz-\frac23y^3\right),\\
 \Delta&=x^2w^2+2wq^j\left(xyz-\frac23y^3\right)
       +\frac19q^{2j}z^2(4xz-3y^2)
\end{align*}
は、量子代数の元をそのPBWベクトル空間で指定する。
これらは $N$ 不変であり、
\[
 x^2\Delta=J^2+\beta q^d
\]
が成り立つ。中心は前節の $D_{m+4}$ 型環 $Z$、
中心原点の極大イデアルは
$\mathfrak m_0=(\mathsf S,\mathsf P,\mathsf R)\subset Z$ である。

\subsection{正規元を反転して得たイデアルのprimitive性}

局所化した環の単純加群を部分環へ制限しても、一般には単純にはならない。
ここでは次の補題を使い、収縮のprimitive性を独立に示す。

\begin{lemma}[単純な正規局所化からのprimitive性]
\label{lem:primitive-contraction-v36}
$A$ を高々可算次元の複素整域、$u\in A$ を非零正規元とする。
$A[u^{-1}]$ が単純環ならば、$A$ はprimitive環である。
また $\{0\}$ は $\operatorname{Spec}A$ の局所閉点である。
\end{lemma}

\begin{proof}
まず $\operatorname{Jac}(A)=0$ を示す。
$a\in\operatorname{Jac}(A)$ が $\mathbb C$ 上超越的なら、
$(1-\lambda a)^{-1}$、$\lambda\in\mathbb C^\times$ は一次独立である。
実際、有限個の一次関係に分母の積を掛けると、相異なる一次因子に関する
部分分数の一意性から全ての係数が零となる。
これは可算次元性に反する。
従って $a$ は代数的である。その最小多項式を $t^kg(t)$、$g(0)\ne0$ と書くと、
$g(a)$ は可逆なので $a^k=0$ となる。整域性から $a=0$ である。

$A$ の任意の非零両側イデアル $I$ は、局所化で全環になるので、
ある $n$ について $u^n\in I$ を満たす。
従って全ての非零素イデアルは $u$ を含む。
もし零イデアルがprimitiveでなければ、全primitive idealが $u$ を含み、
$u\in\operatorname{Jac}(A)=0$ という矛盾を得る。
最後の主張は $D(u)\cap\operatorname{Spec}A=\{0\}$ から従う。
\end{proof}

ここで素イデアルにおける $u^n\in I\Rightarrow u\in I$ は、
$u$ の正規性と $(u)^n=(u^n)$ による。
補題は、局所化後の任意の単純加群の制限が単純になるとは主張しない。

\subsection{可換saturationと星積イデアル}

$C=\mathbb C[x,y,z,w]$ を量子化前の可換多項式環とする。
$N(f)=0$、$X_\alpha(f)=\nu f$ を満たす元については
\[
 f\star a=fa,\qquad a\star f=(1+sN)^\nu(a)f
\]
である。従って $f$ の生成する両側星積イデアルは、ベクトル部分空間として
通常の可換主イデアル $fC$ に一致する。
この主張は複数の同様な生成元にも適用できる。
また $x\star a=xa$ なので、正規元 $x$ による収縮は通常の
可換saturation $I:x^\infty$ と一致する。

特に $\Delta$ および $qJ-\lambda x$ は、この条件を満たす。
後者で二項の $X_\alpha$ 重量が一致することには通常CY値を使っている。
以下のイデアル表示は、量子代数の両側イデアルを指定すると同時に、
$C$ の中のベクトル部分空間を指定する。

\begin{definition}
\label{def:origin-primitives-v36}
$\lambda\in\mathbb C^\times$ に対し
\[
 \mathfrak p_\lambda
   =(\Delta,qJ-\lambda x):x^\infty
\]
と置き、さらに
\[
 \mathfrak p_0=(q,w),\qquad
 \mathfrak p_\infty=(x,H),\qquad
 H=w+\frac{\varepsilon}{4}y^{m-2}z^2
\]
と置く。$H$ の式は $\mathcal S/(x)$ で解釈すれば $N$ 不変である。
また
\[
 \mathfrak n_{z_0,w_0}=(x,y,z-z_0,w-w_0),
 \qquad (z_0,w_0)\in\mathbb C^2
\]
とする。
\end{definition}

$\mathfrak p_0$ の表示では、$q=0$ とした後に $N(w)=0$ となることを使う。
同様に $\mathfrak p_\infty$ では、先に $x=0$ としてから $H$ の正規性を使う。
従っていずれも、表示した通常の可換イデアルと星積イデアルが一致する。

\begin{theorem}[全primitive idealの分類]
\label{thm:all-primitives-v36}
$\mathcal S$ のprimitive idealsは、重複なく次の三種類で尽くされる。
\begin{enumerate}
\item 中心の非原点の閉点 $\chi$ に対する
      $\mathfrak m_\chi\mathcal S$。
      商環は $\mathcal D(\mathbb G_m)$ である。
\item 中心原点上の $\mathfrak p_\lambda$、
      $\lambda\in\mathbb P^1(\mathbb C)=\mathbb C\cup\{\infty\}$。
\item 中心原点上の $\mathfrak n_{z_0,w_0}$、
      $(z_0,w_0)\in\mathbb A^2(\mathbb C)$。
\end{enumerate}
これらは全て完全素イデアルである。
第二種の商環は、$\lambda\ne\infty$ では $x$、
$\lambda=\infty$ では $y$ を反転すると
$\mathcal D(\mathbb G_m)$ になる。
従って第三種を除く全primitive quotientの分数斜体は、
一変数Weyl斜体である。
\end{theorem}

\begin{proof}
高々可算次元の複素代数の単純加群に対するDixmierの議論により、
中心 $Z$ は単純加群上でスカラーとして作用する。
非原点の中心指標では、前節のcentral reductionが単純環なので第一種だけが得られる。
以下では中心指標を原点とする。

$x$ は正規元なので、単純加群上で零として作用するか可逆に作用する。
後者では $x$ を反転した単純加群になり、前節の原点局所化の計算を適用できる。
そのprimitive idealsの逆像は
\[
 \lambda\ne0:\quad(J^2+\beta q^d,qJ-\lambda x),
 \qquad \lambda=0:\quad(q,J)
\]
である。$x$ を反転すると $J^2+\beta q^d=x^2\Delta$ なので、
前者の収縮は定義した $\mathfrak p_\lambda$ である。
後者では $q=0$ の下で $J=x^2w$ となり、
$(q,w)$ は $x$ でsaturatedなので収縮は $\mathfrak p_0$ である。

$x=0$ の場合は、商環で正規な $y$ に同じ議論を適用する。
$y$ が可逆なら、前節のもう一つの局所化は
\[
 \mathbb C[y^{\pm1},H][\tau;sX_\alpha]
\]
であり、中心原点への特殊化が $H=0$ と同値になる。
従って唯一の収縮は $\mathfrak p_\infty=(x,H)$ である。
$y=0$ なら商環は $\mathbb C[z,w]$ であり、第三種だけが得られる。

逆に、$\mathfrak p_\lambda$ の商環は、単純整域
$\mathcal D(\mathbb G_m)$ に埋め込まれ、その正規局所化がこの環となる。
$\lambda=0$ と $\infty$ でも、可換多項式としての表示から各商の
$x$-torsionまたは $y$-torsionがないことが分かり、同じ埋め込みを得る。
従って補題\ref{lem:primitive-contraction-v36}により全てprimitiveである。
第三種の商は $\mathbb C$ なのでprimitiveであり、完全素性も同時に従う。

$\lambda$ の相異なる二点は、局所中心 $c=qJ/x$ の値により区別できる。
$\mathfrak p_\infty$ は $x$ を含み $y$ を含まず、
第三種は両方を含むので、種類の間の重複もない。
最後の分数斜体の主張は、$\mathcal D(\mathbb G_m)$ が
一変数Weyl代数の $\eta$ による局所化であることから従う。
\end{proof}

\begin{corollary}[原点fiberの根基]
\label{cor:origin-radical-v36}
$\mathcal S/\mathfrak m_0\mathcal S$ のprime radicalとJacobson radicalは、
ともに $(\Delta)/\mathfrak m_0\mathcal S$ である。
これは冪零イデアルであり、その商 $\mathcal S/(\Delta)$ は整域である。
\end{corollary}

\begin{proof}
$x$ を反転した後では $\Delta$ は単元倍の
$J^2+\beta q^d$ であり、$d$ が奇数なので後者は既約である。
一方、$\Delta$ の $x=0$ での値は
\[
 \Delta\big|_{x=0}=-\frac43\varepsilon y^dH
\]
という非零多項式であり、$x$ では割れない。
従って $\Delta$ は $C$ でも既約であり、$(\Delta)$ は $x$ でsaturatedである。
星積商は $x$ 局所化のOre整域へ埋め込まれるので、$\mathcal S/(\Delta)$ は整域となる。

中心生成元は $N$ 不変かつ $X_\alpha$ 重量零であるから、
$\mathfrak m_0\mathcal S$ はベクトル部分空間として通常の可換イデアル
$(\mathsf S,\mathsf P,\mathsf R)C$ に一致する。
その零点で $\Delta\ne0$ と仮定すると、
$\mathsf S=q^2\Delta=0$ から $q=0$、
$\mathsf P=2xJ\Delta^L=0$ と $x^2\Delta=J^2$ から $x=J=0$ を得る。
従って $y=0$ であり、$\Delta$ の明示式から $\Delta=0$ となって矛盾する。
可換Nullstellensatzにより、ある $n$ について
$\Delta^n\in\mathfrak m_0\mathcal S$ である。
$\Delta$ は正規元なので、その生成する商イデアルも冪零となる。
整域商を持つことからprime radicalの主張を得る。
また $\mathcal S/(\Delta)$ は可算次元の整域なので、
補題\ref{lem:primitive-contraction-v36}冒頭の議論によりJacobson radicalが零である。
これでJacobson radicalについても同じ等式を得る。
\end{proof}

この系は冪零指数の正確な値を主張しない。
$x$ 局所化後の指数 $L+1$ は、元の原点fiberの指数に対する下界を与える。

\subsection{一つのprimitive idealが含む一次元表現}

原点上の二つの族を
\[
 \mathscr E=\{\mathfrak p_\lambda:\lambda\in\mathbb P^1(\mathbb C)\},\qquad
 \mathscr F=\{\mathfrak n_{z_0,w_0}:(z_0,w_0)\in\mathbb A^2(\mathbb C)\}
\]
と書き、$\mathscr L\subset\mathscr F$ を $w_0=0$ の直線とする。

\begin{lemma}[境界に現れる直線]
\label{lem:primitive-boundary-v36}
全ての $\lambda\in\mathbb P^1(\mathbb C)$ について
\[
 \mathfrak p_\lambda\subseteq\mathfrak n_{z_0,w_0}
 \quad\Longleftrightarrow\quad w_0=0
\]
が成り立つ。相異なる $\mathfrak p_\lambda$ 同士に包含関係はない。
\end{lemma}

\begin{proof}
$\lambda=0$ と $\infty$ は、定義したイデアルに $x=y=0$ を代入すれば直ちに分かる。
以下 $\lambda\ne0,\infty$ とし、通常の可換整域
$A_\lambda=C/\mathfrak p_\lambda$ を使ってsaturationの境界を調べる。
$\mathfrak p_\lambda$ は $x$ 局所化後の可換素イデアルの収縮でもあるから、
$A_\lambda$ は整域である。
その関数体で
\[
 q=-u^2,\quad J=\frac23u^d,\quad
 x=\kappa u^A,\quad y=uv,\qquad
 \kappa=-\frac{2}{3\lambda}
\]
と置ける。実際
$u=-3\varepsilon J/(2q^{j+1})$ とすれば、最初の二式が得られる。
$u^2=-q$ であり、さらに
\[
 v^2=1+\kappa u^dz
\]
なので、$A_\lambda[u,v]$ は $A_\lambda$ の有限整拡大である。
$J$ の式は
\[
 3\kappa^2u^{d+4}w
   =2+\varepsilon(3v-v^3)
   =-\varepsilon(v-2\varepsilon)(v+\varepsilon)^2
\]
を与える。
$x=y=0$ の閉点の上では $u=0$、$v^2=1$ であり、右辺も零なので
$v=-\varepsilon$ でなければならない。
従って $v-\varepsilon$ はそこで単元である。
$(v-\varepsilon)(v+\varepsilon)=\kappa u^dz$ を代入し、整域内で $u$ の冪を消去すると
\[
 \boxed{
 w=-\frac{\varepsilon}{3}\,
       u^{d-4}z^2\frac{v-2\varepsilon}{(v-\varepsilon)^2}.
 }
\]
$d\ge5$ なので $u=0$ では $w=0$ となる。
有限整拡大のlying-overにより、これは $A_\lambda$ 自身の全境界点にも適用される。

逆に
\[
 D_\lambda=
 \frac{\mathbb C[u,v,z,(v-\varepsilon)^{-1}]}
      {(v^2-1-\kappa u^dz)}
\]
の中で、上の式により $x,y,w$ を指定する。
定義方程式は $z$ に関して一次で、係数 $u^d$ と $v^2-1$ が互いに素なので、
Gaussの補題により $D_\lambda$ は整域である。
$u$ を反転した後では $\Delta=0,qJ=\lambda x$ を満たし、
$D_\lambda$ は整域なので、この代入はsaturated ideal
$\mathfrak p_\lambda$ 上でも零となる。
任意の $z_0$ について $u=0,v=-\varepsilon,z=z_0$ と置けば、
$x=y=w=0,z=z_0$ の点を得る。従って境界には直線全体が現れる。

相異なる有限パラメータ間の包含は、局所中心 $c$ の値を比較して排除できる。
$\mathfrak p_\infty$ は $x$ を含むので有限パラメータのイデアルには含まれない。
逆方向については、$\lambda\ne0$ の $qJ-\lambda x$ の像が
$\mathcal S/\mathfrak p_\infty$ で非零の $y^{d+2}=y^A$ の定数倍になることを使う。
$\lambda=0$ では $q$ の像が $-y^2\ne0$ である。
\end{proof}

この証明に使った可換環 $A_\lambda$ は、saturationの包含関係を計算するためのものである。
量子商の積を可換積と同一視しているわけではない。

\subsection{中心原点上の位相の完全な記述}

$\operatorname{Prim}_0(\mathcal S)=V(\mathfrak m_0\mathcal S)$ に
primitive spectrumのJacobson位相を入れる。
定理\ref{thm:all-primitives-v36}により、集合として
$\operatorname{Prim}_0(\mathcal S)=\mathscr E\sqcup\mathscr F$ である。

\begin{theorem}[原点fiber内の任意の集合の閉包]
\label{thm:primitive-closure-v36}
$Y\subseteq\operatorname{Prim}_0(\mathcal S)$ とし、
$Y_E=Y\cap\mathscr E$、$Y_F=Y\cap\mathscr F$ と置く。
$\overline{Y_F}^{\,\mathbb A^2}$ は通常のZariski閉包を表す。
このとき
\[
 \overline Y=
 \begin{cases}
  \operatorname{Prim}_0(\mathcal S),& |Y_E|=\infty,\\[2pt]
  Y_E\cup\mathscr L\cup\overline{Y_F}^{\,\mathbb A^2},
       &0<|Y_E|<\infty,\\[2pt]
  \overline{Y_F}^{\,\mathbb A^2},&Y_E=\varnothing.
 \end{cases}
\]
従って $\mathscr E$ の部分空間位相は
$\mathbb P^1(\mathbb C)$ のZariski位相であり、
$\mathscr F$ は閉部分空間として通常の $\mathbb A^2(\mathbb C)$ である。
各 $\mathfrak p_\lambda$ の一点閉包は
$\{\mathfrak p_\lambda\}\cup\mathscr L$ である。
\end{theorem}

\begin{proof}
$\mathscr F=V(x,y)$ であり、その商環は $\mathbb C[z,w]$ だから、
$\mathscr F$ の位相は主張の通りである。
一点閉包は補題\ref{lem:primitive-boundary-v36}と分類定理から分かる。
有限個の一点と $Y_F$ の和集合については、閉包が有限和と可換であることを使えばよい。

残る場合には、無限個の相異なる $\lambda\in\mathbb C^\times$ が $Y_E$ に含まれる。
それらについて
\[
 \bigcap_{\mathfrak p_\lambda\in Y_E}\mathfrak p_\lambda=(\Delta)
\]
を示す。$x$ を反転し $J^2+\beta q^d$ で割った後、さらに $q$ を反転すると、
係数環は
\[
 \mathbb C[\eta^{\pm1},c^{\pm1}],\qquad
 \eta=J/q^{(d-1)/2},\quad c=qJ/x
\]
となる。$\mathfrak p_\lambda$ はここで $(c-\lambda)$ になる。
非零Laurent多項式が相異なる無限個の $c-\lambda$ で割れることはない。
Ore変数 $t$ の係数ごとにこれを適用すると、交わりは零である。
元の係数環とOre環は整域なので、さらに行った $q$ 局所化は単射である。
従って $x$ 局所化前への収縮は正確に $(\Delta)$ となる。

系\ref{cor:origin-radical-v36}により
$V(\Delta)=\operatorname{Prim}_0(\mathcal S)$ なので、無限集合の場合の主張が従う。
\end{proof}

特にこの空間は $\mathbb P^1$ と $\mathbb A^2$ の位相的な直和ではない。
$\mathscr E$ の各一点は同じ直線 $\mathscr L$ に特殊化する一方、
$\mathscr E$ の任意の無限部分集合は、$w_0\ne0$ の一次元表現にも集積する。
この違いは、パラメータ集合を列挙するだけでは捉えられない。

\subsection{Dixmier--Moeglin同値}

\begin{corollary}
\label{cor:dme-v36}
$\mathcal S$ はDixmier--Moeglin同値を満たす。すなわち素イデアル $P$ について
\[
 P\text{ がprimitive}
 \quad\Longleftrightarrow\quad
 P\text{ がrational}
 \quad\Longleftrightarrow\quad
 \{P\}\text{ が }\operatorname{Spec}\mathcal S\text{ で局所閉}
\]
である。
\end{corollary}

\begin{proof}
可算次元のnoetherian複素代数にはNullstellensatzが適用でき、
一般に「局所閉 $\Rightarrow$ primitive $\Rightarrow$ rational」が成り立つ
\cite{BellWangYee2019}。
本例で残る含意を直接示す。

$P$ がrationalなら、Goldie商環の中心は $\mathbb C$ である。
従って元の中心 $Z$ の像はスカラー環となり、$P$ はある閉点の中心指標を持つ。
非原点ならcentral reductionの単純性から $P=\mathfrak m_\chi\mathcal S$ である。
原点なら系\ref{cor:origin-radical-v36}により $\Delta\in P$ である。
$x\notin P$ の場合、素環 $\mathcal S/P$ 内の非零正規元 $x$ は正則である。
その正規局所化の中心元 $c=qJ/x$ もGoldie商環でスカラーとなる。
$c=\lambda$ の各ファイバーはprime radicalを除くと単純環であるから、
$P=\mathfrak p_\lambda$ を得る。
$x\in P$、$y\notin P$ なら同様に $P=\mathfrak p_\infty$ であり、
$x,y\in P$ なら可換商 $\mathbb C[z,w]$ のrational primeは極大イデアルなので第三種になる。
従ってrational primeは全て分類定理のリストに入る。

第一種と第三種は極大イデアルなので局所閉である。
第二種は補題\ref{lem:primitive-contraction-v36}をその商に適用すれば局所閉となる。
これでrational $\Rightarrow$ 局所閉が従う。
\end{proof}

Lecoutre--Sierraは、線形な局所冪零導分から作る別のCGG族について
Dixmier--Moeglin同値を証明している\cite{LecoutreSierra2019}。
ここでの系は、その一般的現象との整合を本族の明示分類から確認するものである。
本節が新たに与える具体的な情報は、非線形共鳴族の全primitive idealsの表示と、
中心原点上の全ての閉包関係である。
中心指標が非原点を動く無限集合の原点への集積については、ここでは分類していない。

本節は単純加群の零化イデアルを分類しており、単純加群そのものの同型類を分類したものではない。
また通常CY値の次数を忘れた代数を扱っているため、Zhang twistで同値な射影圏の
任意の座標環へprimitive spectrumの結論をそのまま移すことはできない。

\subsection{原点の根基の正確な冪零指数}
\label{subsec:global-origin-nilindex-v37}

系\ref{cor:origin-radical-v36}で得た根基の冪零指数は、正規元を反転した後の指数より一つ大きい。
この差は、境界で失われる情報を測る具体例になる。

\begin{theorem}[全環での原点の厚み]
\label{thm:global-origin-nilindex-v37}
$r=1$、$m=2j+1\ge3$、$s\ne0$ とし、通常CY代表を取る。
$\mathcal R=(\Delta)/\mathfrak m_0\mathcal S$ と置くと
\[
 \boxed{\mathcal R^{\,j+4}=0,\qquad
        \mathcal R^{\,j+3}\ne0.}
\]
従って正確な冪零指数は $(m+7)/2$ である。
$x$ 局所化後の指数 $(m+5)/2$ と比べて一つ増える。
\end{theorem}

\begin{proof}
全ての計算をPBWベクトル空間 $C=\mathbb C[x,y,z,w]$ で行う。
中心生成元が作る星積イデアルは通常の多項式イデアル
$I=(\mathsf S,\mathsf P,\mathsf R)C$ と一致する。
また $\Delta$ は $N$ 不変な正規元なので、星積冪は通常冪と等しく、
根基の $n$ 乗が零であることは $\Delta^n\in I$ と同値である。

\[
 f=xyz-\frac23y^3,\qquad
 g=\frac19z^2(4xz-3y^2),\qquad
 C_j=\frac12(xw^3+q^jyz w^2)
\]
と置き、次の多項式を使う。
\begin{align*}
 T_j={}&2q^{j-1}x^2yw^3
       +q^{2j-2}w^2(2x^2g+4f^2-xyzf)\\
      &+4q^{3j-2}wfg+q^{4j-2}g^2.
\end{align*}
$j\ge1$ なので全ての指数が非負である。
$3f-xyz=2yq$ と
$J=x^2w+q^jf$、$\Delta=x^2w^2+2q^jwf+q^{2j}g$ を用いると
\[
 \Delta^2-2xC_jJ=q^2T_j
\]
を得る。両辺に $\Delta^{j+2}$ を掛ければ
\begin{equation}
 \boxed{\Delta^{j+4}
      =C_j\mathsf P+\mathsf S\Delta^{j+1}T_j.}
 \label{eq:global-nilindex-certificate-v37}
\end{equation}
これが上界を与える。$C_j$ 自身の $N$ 不変性は必要ない。
$\mathsf S,\mathsf P$ は中心元なので、右辺は星積のイデアルにも属する。

下界には、可換多項式環の評価 $y=z=0$ を使う。
この評価は星積環の準同型であるとは主張しない。
上で確認したベクトル空間としてのイデアルの一致に対して適用すると
\[
 I\longmapsto(2x^{2j+7}w^{2j+5})\subset\mathbb C[x,w],
 \qquad
 \Delta^{j+3}\longmapsto x^{2j+6}w^{2j+6}.
\]
後者は前者に属さないので $\Delta^{j+3}\notin I$ である。
\end{proof}

この証明は全ての $j\ge1$ に同一の恒等式を与える。
付属コードでは $q^j$ を独立変数 $t$ に置き換えた恒等式も照合するため、
有限個の重みの計算から指数を外挿しているわけではない。

\section{全正共鳴の素スペクトルと葉空間}
\label{sec:poisson-spectral-v37}

前節までは $r=1$ の中心ファイバーを明示した。本節では $r>0$ の全ての
通常CY代表について、中心を経由せずに素スペクトル全体を記述する。
CGG積を入れる前後で、素イデアルは同じ多項式のベクトル部分空間となる。
線形な局所冪零導分に対する先行結果\cite[Theorem~7.1]{LecoutreSierra2019}と
Ore拡大に対する対応\cite[Theorem~3.6]{JordanOre2014}を踏まえ、
ここでは非線形項 $q^jz^r\partial_w$ に必要な正規局所化を全て明示する。

\subsection{正規元とスライスによる両側イデアルの回収}

\begin{lemma}[非零固有重量を持つスライス]
\label{lem:normal-slice-ideals-v37}
$C$ を有限生成可換整域とし、導分 $N,X$ は
$N$ が局所冪零、$[N,X]=N$ を満たすとする。
$u\in\ker N$、$X(u)=a u$、$a\ne0$ であり、
\[
 C[u^{-1}]=K[t],\qquad K=(\ker N)[u^{-1}],\qquad
 N(t)=1,\quad X(t)=t
\]
と仮定する。$s\ne0$ のCGG積を入れた環を $S^\circ$ とする。
このとき、$S^\circ$ の全両側イデアルと、$K$ の全 $X$ 安定イデアルは
\[
 Q\longmapsto Q[t],\qquad I\longmapsto I\cap K
\]
で対応する。右辺の $Q[t]$ は通常の多項式のベクトル部分空間を意味する。
Poisson括弧 $N\wedge X$ に関する全Poissonイデアルも、同じ対応で与えられる。
両方の対応における素イデアルは、$X$ 安定な通常の素イデアル $Q$ から来る。
従って $S^\circ$ の全素イデアルは完全素である。
\end{lemma}

\begin{proof}
$T=1+sN$ とする。$u\star f=uf$、$f\star u=T^a(f)u$ なので、
局所化内の内部自己同型は $u^{-1}\star f\star u=T^a(f)$ である。
両側イデアル $I$ と $f\in I$ に対し、$T^{an}f\in I$ が全ての整数 $n$ について成り立つ。
$N^{d+1}f=0$ とすれば
\[
 T^{an}f=\sum_{k=0}^{d}\binom{an}{k}s^kN^kf.
\]
$n=0,\ldots,d$ での係数行列は可逆である。
実際、$\binom{an}{k}$ は $n$ の次数 $k$ の多項式で、最高係数 $a^k/k!$ は非零である。
従って $N^kf\in I$、特に $I$ は $N=\partial_t$ に安定となる。

$f=\sum_{i=0}^{d}q_it^i\in I$ と書くと、$N^df/d!=q_d\in I\cap K$ である。
$q_d\star t^d=q_dt^d$ を引き、次数に関して帰納すれば
$I=(I\cap K)[t]$ を得る。
$q\in K$ に対し $[t,q]_\star=sX(q)$ なので、$Q=I\cap K$ は $X$ 安定である。
逆に $Q$ が $X$ 安定なら、$Q[t]$ は $N,X$ の双方に安定であり、
CGG積の各項から両側イデアルとなる。

Poisson側では $\{u,f\}=-auN(f)$ から $N$ 安定性を得る。
同じ次数除去と $\{t,q\}=X(q)$ によって、全く同じイデアル対応が従う。
また、二つの $X$ 安定イデアルについて
\[
 Q_1[t]\star Q_2[t]=(Q_1Q_2)[t]
\]
である。一方の包含は $N,X$ 安定性、他方は
$q_1\star q_2\star t^k=q_1q_2t^k$ から従う。
従って両側素性は $X$ 素性と同値である。
標数零の可換Noether環の微分素イデアルは通常の素イデアルである
\cite[Lemma~2.3]{JordanOre2014}。
最後に $S^\circ/Q[t]\simeq(K/Q)[t;sX]$ は整域上の微分Ore拡大なので整域となる。
\end{proof}

\subsection{全chartの対応と素スペクトルの位相}

以下、$m=2j+r\ge3$、$r>0$ とし、
\[
 \alpha=-\frac{m+r+3}{m+3},\qquad
 C=\mathbb C[x,y,z,w],\quad
 \pi=N\wedge X_\alpha,\quad S=\mathcal S_{j,r}(\alpha,s),\quad s\ne0
\]
と置く。$C$ と $S$ はCGG表示による同じベクトル空間とする。

\begin{theorem}[素スペクトルの同一視]
\label{thm:prime-poisson-correspondence-v37}
$C$ の部分空間 $P$ について、次の二条件は同値である。
\begin{enumerate}
\item $P$ は $S$ の両側素イデアルである。
\item $P$ は $(C,\pi)$ のPoisson素イデアルである。
\end{enumerate}
全ての $S/P$ は整域であり、恒等対応は同相
\[
 \operatorname{Spec}S\simeq\operatorname{PSpec}(C,\pi)
\]
を与える。
\end{theorem}

\begin{proof}
補題\ref{lem:normal-slice-ideals-v37}を次の順序で用いる。
表の $K$ は、指定した閉部分で商を取った後の局所化の $N$ 不変環である。
\[
\begin{array}{c|c|c|c}
 \text{部分}&u& a=X_\alpha(u)/u&t\quad;\quad K\\\hline
 x\ne0&x&\alpha&y/x\ ;\ \mathbb C[x^{\pm1},q,J]\\
 x=0,\ y\ne0&y&-r/(m+3)&z/(2y)\ ;\ \mathbb C[y^{\pm1},H_r]\\
 j=0,\ x=y=0,\ z\ne0&z&3/(m+3)&-w/z^r\ ;\ \mathbb C[z^{\pm1}]
\end{array}
\]
第二行では
\[
 H_r=w+\frac{(-1)^j}{2(r+1)}y^{2j-1}z^{r+1}.
\]
$j=0$ の負の $y$ 冪は $y$ を反転した環の中で用いる。
各行で $N(t)=1$、$X_\alpha(t)=t$、$N(H_r)=0$ が直接確認でき、
通常の可換座標変換により局所化は $K[t]$ となる。
表の固有重量は全て非零である。

$j\ge1$ の場合、残る閉部分 $x=y=0$ では $N=0$ となり、
量子環は可換な $\mathbb C[z,w]$、Poisson括弧は零となる。
$j=0$ の場合、第三行の後に残る閉部分 $x=y=z=0$ は、同様に $\mathbb C[w]$ である。

各段階の $u$ は正規元であり、$(u)$ は通常の多項式イデアルと同じ部分空間である。
また $u^{-n}\star f=u^{-n}f$ なので、局所化と縮約で取る部分空間も一致する。
従って、$u$ を含まない素イデアルを補題で処理し、含む素イデアルを商へ送る帰納によって、
全ての素イデアルの一致と完全素性が従う。
閉集合は任意の部分集合の包含条件 $P\supset E$ で定まるので、
同じ部分空間による対応は両方向に連続である。
\end{proof}

\subsection{原始イデアルと代数的なシンプレクティック葉}

$A=m+r+3$ と置くと、$D=(m+3)X_\alpha$ は整数重量
\[
 D(x,y,z,w)=(-Ax,-ry,(2j+3)z,(m+3r+3)w)
\]
を持つ。$N$ の指数写像による $\mathbb G_a$ 作用と、$D$ による
$\mathbb G_m$ 作用は、二次元の連結可解代数群 $G$ を生成する。
$[D,N]=-(m+3)N$ がその半直積関係を与える。
Poisson階数零の部分を
\[
 F=\begin{cases}V(x,y),&j\ge1,\\V(x,y,z),&j=0\end{cases}
\]
とする。

\begin{theorem}[葉、原始スペクトル、Dixmier--Moeglin同値]
\label{thm:primitive-leaf-space-v37}
$\mathbb A^4\setminus F$ におけるシンプレクティック葉は、二次元の $G$ 軌道である。
$F$ 上の葉は個々の点である。
葉 $L$ の閉包の消滅イデアルを $I(\overline L)$ とすると、
\[
 \operatorname{Prim}S
 =\operatorname{PPrim}(C,\pi)
 =\{I(\overline L):L\text{ はシンプレクティック葉}\}
\]
が同じ多項式部分空間として成立する。
$S$ はDixmier--Moeglin同値を満たし、$(C,\pi)$ はPoisson版の同値を満たす。
正次元の葉に対応する原始商は、上の表で最初に反転可能な正規元を反転すると
$\mathcal D(\mathbb G_m)$ となる。
\end{theorem}

\begin{proof}
$x\ne0$ では $N(x)=0,N(y)=x$、$X_\alpha(x)\ne0$ なので $N,X_\alpha$ は独立である。
$x=0,y\ne0$ では $N(y)=0,N(z)=2y$、$X_\alpha(y)\ne0$ を用いる。
$j=0,x=y=0,z\ne0$ では $N(z)=0,N(w)=-z^r$、$X_\alpha(z)\ne0$ を用いる。
残る $F$ では $N=0$ である。従って表示したものが階数零部分の全体である。
階数二の部分ではHamilton分布は $N,X_\alpha$ が張る分布であり、
連結な $G$ 軌道と一致する。代数群軌道は局所閉なので、全ての葉は代数的である。

各局所chartは $K[t]$ であり、$G$ 軌道は $\mathbb G_m$ の $\operatorname{Spec}K$ 上の
軌道と $t$ 方向の直積として扱える。
$K$ には非零重量の可逆元 $u$ があるため、その $\mathbb G_m$ 軌道は閉じている。
実際、$u(\lambda\cdot p)=\lambda^a u(p)$、$u(p)\ne0$ なので、軌道写像の座標環の像は
$\lambda^a$ とその逆を含み、
$\mathbb C[\lambda^{\pm1}]$ はその像の上で有限整となる。
従って軌道写像は閉写像となる。
各軌道は $\mathbb G_m$ を有限部分群で割ったもので、再び $\mathbb G_m$ と同型である。
その消滅イデアル $Q\subset K$ で割ると、局所量子商は
\[
 \mathbb C[v^{\pm1}][t;c\,v\partial_v]\simeq\mathcal D(\mathbb G_m),\qquad c\ne0
\]
となる。これは単純環である。
補題\ref{lem:primitive-contraction-v36}により、縮約した素イデアルはprimitiveかつ局所閉である。
$F$ の点は可換な商環の極大イデアルなので、同じ結論を満たす。

逆に、任意の量子rational素イデアル $P$ を取る。
$P$ が $F$ 上にあれば、その可換商の有理性から点のイデアルである。
そうでなければ、該当chartで $P=Q[t]$ とし、$K_0=\operatorname{Frac}(K/Q)$ と置く。
$Q$ は整数重量導分 $D$ に安定なので、有限個の重量成分への補間によって
$\mathbb G_m$ 安定となる。従って $K_0$ は固有ベクトルで生成される。
非零整数重量 $a$ の可逆元 $u$ と、重量 $b$ の固有元 $f$ について
$f^a/u^b$ は $D$ 不変である。負の指数を許す関数体の中で考えればよい。
これにより $K_0$ は $K_0^D(u)$ 上代数的であり、$u$ は $K_0^D$ 上超越的なので
\[
 \operatorname{trdeg}_{\mathbb C}K_0^D
 =\operatorname{trdeg}_{\mathbb C}K_0-1.
\]
$K_0^D$ は量子商の分数斜体の中心に入る。
従ってrational性から $\dim K/Q=1$ となる。
この既約アフィン曲線上では、可逆な非零重量元があるため固定点が存在しない。
全ての軌道は一次元の開集合であり、既約性によって軌道は一つしかない。
ゆえに前段の $\mathbb G_m$ 軌道の場合へ帰着し、$P$ は局所閉かつprimitiveである。
有限生成複素Noether代数に対する通常の含意
「局所閉 $\Rightarrow$ primitive $\Rightarrow$ rational」
\cite[\S8]{LecoutreSierra2019}と合わせてDixmier--Moeglin同値を得る。

Poisson coreについては、点を含む葉の閉包のイデアルが、その点の極大イデアルに含まれる
最大のPoissonイデアルである。
これは各chartでHamilton導分から $N,X_\alpha$ を回収できること、
および葉が代数的な軌道であることから従う。
従って上で列挙したイデアルがPoisson primitiveの全体でもある。
Poisson rational素イデアルにも $K_0^D$ による同じ次元評価を適用できる。
局所閉性は定理\ref{thm:prime-poisson-correspondence-v37}により一致するので、
列挙した正次元の葉のPoisson関数体は $\mathbb C(v,t)$、$\{t,v\}=cv$ ($c\ne0$) で、中心は $\mathbb C$ である。
点の葉についても同じ有理性が成り立つので、Poisson版のDixmier--Moeglin同値も成立する。
\end{proof}

この結果は、局所化した環の単純加群を元の環へ制限すれば単純である、
という一般には誤った推論を使わない。使うのは単純な正規局所化からの
零化イデアルのprimitive性である。

\begin{corollary}[全正共鳴の原始スペクトルの層別]
\label{cor:primitive-strata-v37}
$g=\gcd(m+r+3,r)$、$n=(m+r+3)/g$ とする。
原始スペクトルは次の局所閉部分に分かれる。各部分の内部の位相も表示した多様体の位相と一致する。
\[
\begin{array}{c|c}
 \text{条件}&\text{パラメータ空間}\\\hline
 x\notin P&\{ab=c^n\}(\mathbb C)\\
 x\in P,\ y\notin P&\mathbb A^1(\mathbb C)\\
 j=0,\ x,y\in P,\ z\notin P&\{\mathrm{pt}\}\\
 P\text{ は }F\text{ 上の文字表現}&F(\mathbb C)
\end{array}
\]
第三行は $j=0$ の場合だけ現れる。
この表の直和に離散的な位相を入れるのではなく、異なる部分の貼合せは次の閉包公式で定める。
\end{corollary}

\begin{proof}
第一行は $\mathbb C[x^{\pm1},q,J]$ の $\mathbb G_m$ 商である。
$x=1$ と規格化すると残る有限作用は
$\mu_A$ の重量 $(2r,rd)$ の作用であり、零重量環は
定理\ref{thm:localized-center-v35}の $\mathbb C[a,b,c]/(ab-c^n)$ となる。
第二行では $B=m+3r+3$、$g'=\gcd(r,B)$ とすれば、零重量環は
\[
 \mathbb C[y^{B/g'}H_r^{r/g'}]
\]
である。第三行の零重量環は $\mathbb C$ である。
いずれも、可逆元を一へ規格化した後の有限群商として、閉軌道と零重量環の極大点が対応する。
有限群商の閉集合の対応から各部分の位相の主張も従う。
最後の行は零Poisson構造の可換商である。
\end{proof}

\subsection{中心指標を動かす集合も含む閉包公式}

点 $p$ の葉を $L_p$ とし、
\[
 \kappa:\mathbb A^4(\mathbb C)\longrightarrow\operatorname{Prim}S,
 \qquad p\longmapsto I(\overline{L_p})
\]
と置く。

\begin{theorem}[全域的な閉包位相]
\label{thm:global-leaf-closure-v37}
$\kappa$ は、Zariski位相から原始スペクトルの位相への商写像である。
任意の部分集合 $T\subset\mathbb A^4(\mathbb C)$ について
\[
 \operatorname{Sat}_\pi(T)=G\cdot(T\setminus F)\ \cup\ (T\cap F)
\]
と定めると
\[
 \boxed{\quad
 \overline{\kappa(T)}^{\operatorname{Prim}S}
 =\kappa\!\left(\overline{\operatorname{Sat}_\pi(T)}^{\mathrm{Zar}}\right).
 \quad}
\]
従って、異なる中心指標に属する集合の集積も、同じ代数群軌道の閉包計算で扱える。
\end{theorem}

\begin{proof}
$\kappa$ の全射性は前定理による。
任意の集合 $E\subset C$ に対して
$\kappa^{-1}\{P:E\subset P\}$ は、$E$ の全Hamilton反復微分の共通零点であり、閉じている。
従って $\kappa$ は連続である。
葉の和である閉集合 $V$ の消滅イデアルはPoisson安定なので、
\[
 \kappa(V)=\{P\in\operatorname{Prim}S:I(V)\subset P\}
\]
は閉集合となる。これにより商写像性を得る。

葉の和のZariski閉包も葉の和である。
実際、その消滅イデアルは各葉に接する全Hamilton導分で安定である。
従って $W=\overline{\operatorname{Sat}_\pi(T)}$ は、$T$ を含む最小の閉じた葉の和である。
$\kappa(W)$ は閉集合で $\kappa(T)$ を含み、
逆に $\kappa(T)$ を含む任意の閉集合の逆像は $W$ を含む。
これが表示した等式である。
\end{proof}

この公式は、$r=1$ の原点ファイバーの閉包公式を含む全域的な幾何記述である。
一方、中心曲面と $\mathbb P^1$ のパラメータだけを使った全ての閉集合の一覧や、
部分解消上の代数の層の構成を主張するものではない。
また、$r=0$ では第二chartの固有重量が零になるため、本節の証明をそのまま適用しない。

\section{端点の一の冪根ファイバー：中心、Azumaya部分と二つの単純加群}
\label{sec:endpoint-root-unity-v38}

ここでは正共鳴族から、偶数端点の全変形を実現する明示Ore族へ移る。
一の冪根において、この族は中心上有限な非可換幾何のchartを持つ。
その中心と単純加群は混合項によらず、冪根の位数の偶奇で決まる。
以下の結果は中心の一つの開部分に関する完全な計算であり、
閉境界上の全中心の表示は含まない。

\subsection{混合項を消去できる中心局所化}

$s\gamma\ne0$、$1-s^2\gamma\ne0$、$\rho^2=\gamma$ とし、
\[
 t=\frac{1-s\rho}{1+s\rho},\qquad \beta=\frac1{4\gamma}
\]
と置く。三変数AS正則代数 $B$ を、次数一の $x,u,v$ と関係
\begin{equation}
\label{eq:root-B-v38}
 ux=txu,\qquad vx=t^{-1}xv,\qquad
 uv-vu=(t^{-1}-t)\beta x^2
\end{equation}
で定める。中心元と正規化自己同型は
\[
 Q=-uv-t\beta x^2,\qquad
 \sigma(x)=x,\quad \sigma(u)=tu,\quad \sigma(v)=t^{-1}v.
\]
公表稿Xで構成した座標環は、任意の斉次混合項 $H\in B_{m+1}$ に対して
\[
 S_H=B[w;\sigma,\delta_H],\qquad
 \delta_H(b)=x^{-1}[b,H],\qquad J=xw+H
\]
である。$B/(x)$ の可換性により左割算は $B$ 内で定義され、$J$ は中心元である。
この構成とAS正則性・CY性の証明は公表稿X\cite{VKCategories36}による。
ここで必要な一般化Weyl代数の表示は
\[
 B[x^{-1}]=R(\tau,a),\qquad
 R=\mathbb C[Q,x,x^{-1}],\quad \tau(x)=tx,\quad \tau(Q)=Q,
 \quad a=-Q-t^{-1}\beta x^2,
\]
すなわち $ub=\tau(b)u$、$vb=\tau^{-1}(b)v$、$vu=a$、$uv=\tau(a)$ である。
これは標準的な一般化Weyl代数の規約\cite{SuarezVivas2015Root38}である。

\begin{lemma}[中心の開部分から混合項を消す]
\label{lem:root-localization-v38}
$t$ の位数を有限な $\ell\ge3$ とする。
$\Xi=x^\ell$ は全代数 $S_H$ の中心に入り、
\begin{equation}
\label{eq:root-central-chart-v38}
 A_H:=S_H[\Xi^{-1}]
 \simeq B[x^{-1}][J]
\end{equation}
となる。右辺の $J$ は独立な中心変数である。
特にこの中心局所化は、次数と混合項 $H$ に依存しない非次数付き代数である。
\end{lemma}

\begin{proof}
$\Xi\in Z(B)$ であり、$\sigma(\Xi)=\Xi$、$\delta_H(\Xi)=0$ なので
$w\Xi=\Xi w$ である。
$w\mapsto x^{-1}(J-H)$ により $S_H$ を $B[x^{-1}][J]$ に埋め込める。
単射性は、Ore多項式の最高 $w$ 次係数が最高 $J$ 次係数として残ることから分かる。
$\Xi$ を可逆にすれば $x^{-1}=x^{\ell-1}\Xi^{-1}$ が得られ、
逆に $J=xw+H$ が得られるため、局所化は右辺に一致する。
ここでは $x$ 自身が $S_H$ の正規元であるとは仮定していない。
$t=1$ は $s\rho=0$ を要し、$t=-1$ は表示した分数では生じないので、
有限位数は三以上である。
\end{proof}

\subsection{中心の偶奇公式}

\begin{theorem}[一の冪根chartの中心]
\label{thm:root-center-v38}
$U=u^\ell$、$V=v^\ell$ とする。
$A_H$ は中心上有限でPI次数は $\ell$ である。
その中心は、$\ell$ が奇数なら
\begin{equation}
\label{eq:root-center-odd-v38}
 Z(A_H)=
 \frac{\mathbb C[Q,J,\Xi,\Xi^{-1},U,V]}
 {(UV+\beta^\ell\Xi^2+Q^\ell)},
\end{equation}
$\ell=2d$ が偶数なら
\begin{equation}
\label{eq:root-center-even-v38}
 Z(A_H)=
 \frac{\mathbb C[Q,J,\Xi,\Xi^{-1},U,V]}
 {(UV-(Q^d+(-1)^d\beta^d\Xi)^2)}
\end{equation}
である。また $Z(A_H)=Z(S_H)[\Xi^{-1}]$ である。
\end{theorem}

\begin{proof}
一般化Weyl代数の $\mathbb Z$ 次数を $\deg u=1$、$\deg v=-1$ とする。
$x$ と可換な成分の次数は $\ell$ の倍数であり、
その係数が $u$ とも可換である条件は $\tau$ 不変性である。
従って中心は $R^\tau[J]=\mathbb C[Q,J,\Xi,\Xi^{-1}]$ と $U,V$ で生成される。
関係式は
\begin{equation}
\label{eq:root-norm-v38}
 UV=\prod_{i=0}^{\ell-1}(-Q-t^{2i+1}\beta x^2)
\end{equation}
である。$e=\gcd(2,\ell)$、$d_0=\ell/e$ と置くと、右辺は
\[
 (-1)^\ell\bigl(Q^{d_0}-(-t)^{d_0}\beta^{d_0}x^{2d_0}\bigr)^e.
\]
奇数の場合は $(-t)^{d_0}=-1$、偶数の場合は $(-t)^{d_0}=(-1)^{d_0+1}$ なので、二つの式を得る。
一般化Weyl代数の直和分解
\[
 R[J]\ \oplus\,\bigoplus_{k\ge1}R[J]u^k
 \ \oplus\,\bigoplus_{k\ge1}R[J]v^k
\]
により、表示した関係以外の関係はない。

$0\le i<\ell$ と $0\le k<\ell$ に対する $x^iu^k$、
および $0\le i<\ell$ と $1\le k<\ell$ に対する $x^iv^k$ が
中心上の有限生成系となる。
中心の分数体上では $u$ も可逆であり、$x^i u^k$、$0\le i,k<\ell$ が基底である。
関係 $ux=txu$ により中心上の次元は $\ell^2$、PI次数は $\ell$ となる。

最後の中心の局所化について、$b\Xi^{-n}$ が局所化後に中心的なら、
$[b,c]\Xi^{-n}=0$ が全ての $c\in S_H$ に対して成り立つ。
整域性により $[b,c]=0$ であるから $b\in Z(S_H)$ となる。
\end{proof}

\subsection{Azumaya部分と特異中心ファイバー}

$\mathcal R$ を、二頂点 $e_0,e_1$ の間に各向き二本ずつの矢を持ち、
長さ二の全ての道を零とする六次元代数とする。
従って
\[
 \operatorname{rad}\mathcal R\simeq\mathbb C^4,\qquad
 (\operatorname{rad}\mathcal R)^2=0,\qquad
 \mathcal R/\operatorname{rad}\mathcal R\simeq\mathbb C\oplus\mathbb C.
\]

\begin{theorem}[中心特異点と単純加群の分裂]
\label{thm:root-Azumaya-v38}
中心の任意の複素閉点 $p$ を取る。
\begin{enumerate}
\item $\ell$ が奇数なら中心は滑らかであり、$A_H$ は中心全体上で
次数 $\ell$ のAzumaya代数である。
各中心ファイバーは $M_\ell(\mathbb C)$ である。
\item $\ell=2d$ が偶数なら、中心の特異部分は
\[
 \Sigma=V(U,V,Q^d+(-1)^d\beta^d\Xi)
 \simeq\mathbb G_m\times\mathbb A^1_J
\]
である。Azumaya部分は正確に $\operatorname{Spec}Z(A_H)\setminus\Sigma$ であり、
その各閉点の中心ファイバーは $M_{2d}(\mathbb C)$ である。
\item 偶数の場合、$p\in\Sigma$ では非標準的な代数同型
\begin{equation}
\label{eq:root-singular-fiber-v38}
 A_H\otimes_{Z(A_H)}\mathbb C(p)\simeq M_d(\mathcal R)
\end{equation}
が存在する。従ってこのファイバーは次元 $6d^2$、根基の次元は $4d^2$、
根基の平方は零である。
この中心指標を持つ単純加群は同型を除いて二つであり、どちらも次元 $d$ である。
\end{enumerate}
\end{theorem}

\begin{proof}
奇数の場合、中心方程式の $\Xi$ 微分は $2\beta^\ell\Xi$ であり、
$\Xi\ne0$ なので特異点はない。
偶数の場合、$F=Q^d+(-1)^d\beta^d\Xi$ を新しい座標に取れば
中心は $UV=F^2$ と $F-Q^d\ne0$ で記述される。
従って特異部分は $U=V=F=0$ で、ここでは $Q\ne0$ である。
横断方向の特異点は $A_1$ 型である。

表現とAzumaya性を、一般的な一致定理を使わずに確認する。
中心上の有限エタール拡大 $\lambda^\ell=\Xi$ を取り、
$x$ の相異なる固有値 $\lambda,t\lambda,\ldots,t^{\ell-1}\lambda$ に対応する
冪等元を $e_i$ とする。
$u,v$ はこれらの間を逆向きに結ぶ周期的な矢となり、
同じ辺の二方向の合成は、添字の巡回的な変更を除いて
\[
 a_i=-Q-t^{2i-1}\beta\lambda^2
\]
というスカラーである。
全周を回る二つの道は、それぞれ中心元 $U,V$ である。
$a_i$ が可逆なら、その辺の二本の矢はスカラー倍の互いの逆であり、
その辺を縮約して二頂点の行列単位を同一視できる。

$U\ne0$ または $V\ne0$ の開部分では一方向の全ての矢が可逆となるため、
エタール局所的に全行列代数となる。
$U=V=0$ の点では少なくとも一つの $a_i$ が零である。
$t^2$ の位数は奇数の場合 $\ell$、偶数の場合 $d$ であるから、
零となる辺の数は奇数で一、偶数で二である。
偶数の場合の二辺は周期上でちょうど $d$ 個離れている。

零辺が一つの場合、残る $\ell-1$ 本の可逆辺が全頂点を結ぶ。
それらの辺を縮約すれば、残る二つの矢は中心元 $U,V$ の単元倍となり、
唯一の関係は $UV=\prod a_i$ である。
従って、その点の近傍でエタール局所的に
$M_\ell(Z(A_H))$ となる。
これはファイバーの単純性だけでなく、中心上の局所自由性も証明している。

偶数で $p\in\Sigma$ の場合は、二つの零辺を除くと、各 $d$ 頂点の二本の鎖が残る。
各鎖の可逆辺を縮約して一頂点にすると、各向き二本、合計四本の矢が残る。
同一辺の往復は $a_i=0$ により零である。
異なる二辺を同じ向きに通る道は $U=0$ または $V=0$ により零である。
これで長さ二の道は全て零となる。
この簡約から、元のファイバーの次元は高々 $6d^2$ である。
追加の一次関係がないことは逆向きの表現で確認できる。
二つの鎖の内部では、$u$ を行列単位、逆向きの $v$ を対応する非零 $a_i$ 倍の行列単位とし、
二つの切断辺では $\mathcal R$ の四つの独立な根基矢を配置する。
$x$ は固有値 $\lambda t^i$ の対角行列で表す。
切断辺での往復と全周の道は根基の平方零性により零であり、
その他の関係も $a_i$ の定義により成立する。
この表現の像は両方の $M_d(\mathbb C)$ ブロックと四つの根基ブロックを生成するため、
$M_d(\mathcal R)$ への全射となる。
次元上界と合わせて式\eqref{eq:root-singular-fiber-v38}を得る。
そのファイバーは全行列代数でないためAzumayaではない。
上で調べた開部分と合わせてAzumaya部分の正確な記述となる。
単純加群の数と次元、および根基の記述は $M_d(\mathcal R)$ から直ちに従う。
\end{proof}

\begin{remark}[全代数とchartの区別]
定理は全代数の中心の開部分 $\Xi\ne0$ を決定する。
この開部分での中心、Azumaya部分、単純加群は $H$ に依存しないが、
境界 $\Xi=0$ にも同じ主張が成立するとは限らない。
全次数付き環の同型分類、$\operatorname{qgr}$ の同値分類、
および導来同値分類も、この非次数付き中心局所化の計算とは別の問題である。
\end{remark}

\subsection{検証可能な全位数の仕組み}

中心の式は一つの巡回ノルム\eqref{eq:root-norm-v38}の因数分解から、
表現の分裂は $t^2$ の周期から得られる。
従って有限個の位数の計算から全位数へ外挿していない。
付属検証では円分多項式による厳密な剰余計算でノルムを照合し、
零辺を持つ巡回quiverの中心ファイバーを実際の有限次元代数として検査する。
偶数位数では、一般の中心指標には $2d$ 次元の単純加群があり、
特異中心指標には二つの $d$ 次元単純加群と、平方零の四つの矢が現れる。
この表現論的な厚みは、中心の方程式 $UV=F^2$ のみよりも多くの情報を持つ。

\section{偶数冪根chartの商スタック表示と形式McKay模型}
\label{sec:root-McKay-v38}

偶数位数で現れた二つの単純加群と四つの矢は、
特異中心ファイバーだけの偶然ではない。
有限エタール基底変換の後では、中心の開部分全体が、
滑らかな四次元多様体への二元群の作用によって記述できる。
さらに、各特異点の完成代数は標準的な $A_1$ 型McKay代数の行列環になる。

\subsection{有限エタール基底変換後のMorita表示}

$\ell=2d$ とし、前節の $A_H=S_H[\Xi^{-1}]$ と $Z=Z(A_H)$ を用いる。
\[
 Z'=Z[\lambda]/(\lambda^{2d}-\Xi),\qquad
 A'=A_H\otimes_Z Z'
\]
と置く。$\Xi$ が可逆で特性零なので、$Z'/Z$ は有限エタールである。
以下の可換環を定める：
\begin{equation}
\label{eq:root-smooth-cover-v38}
 C_0=\mathbb C[Q,J,\lambda,\lambda^{-1}],\qquad
 f=(-1)^dQ^d+\beta^d\lambda^{2d},\qquad
 T=\frac{C_0[P,R]}{(PR-f)}.
\end{equation}
二元群 $G=\{1,g\}$ は $C_0$ を固定し、$g(P)=-P$、$g(R)=-R$ で作用する。
$T\#G$ は対応するskew group algebraであり、
$g^2=1$、$gP=-Pg$、$gR=-Rg$ を満たす。

\begin{theorem}[滑らかな二元群商による記述]
\label{thm:root-etale-Morita-v38}
$T$ は滑らかな四次元整域であり、
\[
 T^G\simeq Z',\qquad
 A'\ \text{は}\ T\#G\ \text{とMorita同値}
\]
である。従って
\[
 \operatorname{Perf}(A')\simeq
 \operatorname{Perf}\bigl([\operatorname{Spec}T/G]\bigr).
\]
これは $\Xi\ne0$ における有限エタール基底変換後の記述であり、
元の $A_H$ が大域的に行列環になるという主張ではない。
\end{theorem}

\begin{proof}
$x$ の固有値 $\lambda t^i$ に対応する冪等元を
\[
 e_i=\frac1{2d}\sum_{k=0}^{2d-1}t^{-ik}(x/\lambda)^k
 \qquad(0\le i<2d)
\]
とし、$e=e_0+e_d$ と置く。
角代数 $eA'e$ の中で
\[
 g=e_0-e_d,\qquad P=e u^d e,\qquad R=e v^d e
\]
と置けば、$xe=\lambda g$ であり、
$g^2=e$、$gP=-Pg$、$gR=-Rg$ となる。
半周のノルムは
\[
 u^dv^d=v^du^d
 =\prod_{i=0}^{d-1}(-Q-t^{2i+1}\beta x^2)
 =(-1)^dQ^d+\beta^d\Xi=f.
\]
従って $PR=RP=fe$ であり、$T\#G\to eA'e$ という準同型を得る。

この写像が同型であることを、一般化Weyl代数の次数分解で確かめる。
$u^n$ または $v^n$ が $e_0,e_d$ の間を結ぶための条件は $d\mid n$ である。
また零次の角は $C_0e_0\oplus C_0e_d$ である。
従って $eA'e$ は $C_0$ 加群として
\[
 (C_0\oplus C_0g)
 \ \oplus\,\bigoplus_{k\ge1}(C_0\oplus C_0g)P^k
 \ \oplus\,\bigoplus_{k\ge1}(C_0\oplus C_0g)R^k
\]
という直和分解を持つ。
$PR=f$ を用いた $T\#G$ の正規形も正確にこの直和分解であるから、
準同型は単射かつ全射である。
ここで $A'$ の整域性は使っていない。

次に $e$ がfullであることを示す。
中心の任意の閉点で、$U\ne0$ または $V\ne0$ なら、
巡回quiverの全頂点は可逆矢で結ばれ、$e$ はfullになる。
$U=V=0$ なら二つの零辺が距離 $d$ だけ離れており、
残る二つの鎖はいずれも $d$ 頂点を持つ。
互いに反対の頂点 $0,d$ は異なる鎖に入るので、$e$ は各鎖を一つずつ代表する。
可逆辺に沿う行列単位により、この場合も各ファイバーでfullである。
$A'/A'eA'$ は中心上有限な加群で、全ての閉点ファイバーが零である。
$Z'$ が有限型の複素代数であることとNakayamaの補題から、
$A'/A'eA'=0$ となる。よって $A'$ と $eA'e=T\#G$ はMorita同値である。

関係式 $PR-f$ の $\lambda$ 微分は
$-2d\beta^d\lambda^{2d-1}$ であり、$T$ の単元である。
従って $T$ は滑らかである。$PR-f$ の標準形から整域性も分かる。
不変部分は
\[
 T^G\simeq\frac{C_0[U,V]}{(UV-f^2)}
 \simeq
 \frac{\mathbb C[Q,J,\lambda,\lambda^{-1},U,V]}{(UV-f^2)}=Z'.
\]
ここで $U=P^2$、$V=R^2$ であり、
$f=(-1)^d(Q^d+(-1)^d\beta^d\Xi)$ なので、前節の中心式と一致する。
最後に、$T\#G$ 加群と $G$ 同変な $T$ 加群の同値を完全対象へ制限すれば、
表示した商スタックとの同値を得る。
\end{proof}

\subsection{特異中心の完成代数}

\begin{corollary}[全ての偶数位数に共通する形式模型]
\label{cor:root-formal-McKay-v38}
$p\in\Sigma$ とし、中心に関する完成を
\[
 \widehat Z_p=\widehat{Z_p},\qquad
 \widehat A_p=A_H\otimes_Z\widehat Z_p
\]
と書く。形式変数 $a,b,z_1,z_2$ を適切に選べば、
\begin{equation}
\label{eq:root-formal-McKay-v38}
 \widehat A_p\simeq
 M_d\bigl(\mathbb C[[a,b,z_1,z_2]]\#\mu_2\bigr),
\end{equation}
ここで非自明な元は $(a,b,z_1,z_2)\mapsto(-a,-b,z_1,z_2)$ と作用する。
中心の同型は
\[
 \widehat Z_p\simeq
 \frac{\mathbb C[[A,B,F,z_1,z_2]]}{(AB-F^2)},
 \qquad A=a^2,\quad B=b^2,\quad F=ab
\]
である。この完成代数は、その完成中心の非可換クレパント解消である。
\end{corollary}

\begin{proof}
残差体が $\mathbb C$ で $\Xi(p)\ne0$ なので、
$\widehat Z_p$ 内で $\lambda^{2d}=\Xi$ の根を一つ選べる。
これはHenselの補題による。従って先の有限エタール表示は、
完成中心上では一つの成分を選ぶだけで利用できる。
二つの零辺以外は完成環上で可逆である。
添字を巡回的に付け直し、各 $d$ 頂点の鎖の内部の行列単位を使うと、
\[
 \widehat A_p\simeq M_d(e\widehat A_pe)
\]
を得る。二つの鎖が同じ長さを持つことが、この行列サイズを決めている。

特異点では $P=R=0$、$Q=Q_0\ne0$、$\lambda=\lambda_0\ne0$ である。
\[
 \frac{\partial f}{\partial Q}(p)=(-1)^d dQ_0^{d-1}\ne0
\]
なので、形式的陰関数定理により $PR=f$ から $Q$ を消去できる。
$a=P$、$b=R$、$z_1=\lambda-\lambda_0$、$z_2=J-J(p)$ と取ると、
$T$ の該当する完成は $\mathbb C[[a,b,z_1,z_2]]$ に等しい。
$Q$ の解は $ab,\lambda,J$ の関数であるため、この同一視は $\mu_2$ 同変である。
これが式\eqref{eq:root-formal-McKay-v38}を与える。
不変環の式は、偶数の単項式が $a^2,ab,b^2,z_1,z_2$ で生成されることによる。

最後の主張は、商特異点に対する標準的な非可換クレパント解消
\cite[Example~1.1]{VanDenBerghMcKay38}の、この場合の完成版である。
具体的に $L=\mathbb C[[a,b,z_1,z_2]]$ と置くと、
$\mu_2$ の固定部分は余次元二であり、余次元一では作用は自由である。
従って自然写像 $L\#\mu_2\to\operatorname{End}_{L^{\mu_2}}(L)$ は余次元一で同型である。
両辺は正規環 $L^{\mu_2}$ 上の反射的加群なので、全体で同型となる。
$L$ と $L\#\mu_2$ は不変環上で最大Cohen--Macaulayである。
また群の位数が可逆なので、$L$ 上の射影加群に対する分裂を平均化でき、
$\operatorname{gldim}(L\#\mu_2)=\operatorname{gldim}L=4$ となる。
例えば上界は $L$ 上の四段目のsyzygyの射影性と平均化から、
下界は残差体の射影次元を $L$ へ制限することで得られる。
従って $L\#\mu_2$ は不変環の非可換クレパント解消であり、
その行列環は反射的加群 $L^{\oplus d}$ の自己準同型環として同じ性質を持つ。
\end{proof}

中心の閉点への特殊化では $a^2=ab=b^2=0$ となるため、
\[
 \bigl(\mathbb C[a,b]/(a,b)^2\bigr)\#\mu_2
\]
が現れる。二つの群指標が二頂点を与え、$a,b$ が各方向二本の矢を与える。
これは前節の六次元代数 $\mathcal R$ そのものである。
従って平方零の四矢は、全形式近傍のMcKay構造からも説明できる。
この形式Morita型は $d$ に依存せず、位数による違いは行列サイズ $d$ に残る。
ここでの非可換クレパント解消という結論は完成中心についてのものであり、
元の中心上で大域的な自己準同型環表示が得られたとは主張しない。

\section{特異部分への降下、単純加群のモノドロミーとBrauer障害}
\label{sec:singular-descent-v39}

前節の有限エタール表示は、各特異中心指標に二つの単純加群があることを説明した。
しかし、この二つを中心の特異部分全体にわたって別々に番号付けできるとは限らない。
本節では元の中心上に戻り、特異部分に制限した代数そのものを決定する。
二つの単純加群は連結な二重被覆上で初めて分離する。
さらに、滑らかな中心部分にあるAzumaya代数のBrauer類は正確に位数 $\ell$ を持つ。
この類は、完成局所的な非可換クレパント解消を元の中心上の自己準同型環表示へ降下することを妨げる。

\subsection{特異部分上の代数全体}

まず $\ell=2d\ge4$ とし、前節の記号
\[
 A=A_H=B[x^{-1}][J],\qquad Z=Z(A),\qquad
 F=Q^d+(-1)^d\beta^d\Xi
\]
を用いる。特異部分は
\[
 \Sigma=V(U,V,F)\subset\operatorname{Spec}Z,
 \qquad R_\Sigma=\mathbb C[Q,Q^{-1},J],
 \qquad A_\Sigma=A\otimes_ZR_\Sigma
\]
であり、ここでは $\Xi=(-1)^{d+1}(Q/\beta)^d$ である。
二次エタール代数とその非自明な自己同型を
\begin{equation}
\label{eq:singular-double-cover-v39}
 L_\Sigma=R_\Sigma[\kappa]/(\kappa^2+tQ/\beta),
 \qquad \iota(\kappa)=-\kappa
\end{equation}
とする。$L_{\Sigma,\iota}$ は左作用が通常の乗法、右作用が $\iota$ で捻られた
$L_\Sigma$ 双加群を表す。

\begin{theorem}[特異部分上の大域的な降下代数]
\label{thm:singular-descent-v39}
$R_\Sigma$ 代数として
\begin{equation}
\label{eq:singular-descent-algebra-v39}
 A_\Sigma\simeq M_d(\mathscr C_\Sigma),\qquad
 \mathscr C_\Sigma=
 L_\Sigma\ltimes
 \bigl(L_{\Sigma,\iota}\oplus L_{\Sigma,\iota}\bigr)
\end{equation}
となる。右辺の付加部分は平方零である。
具体的に $\mathscr C_\Sigma$ は $\kappa,p,r$ で生成され、関係は
\begin{equation}
\label{eq:singular-corner-relations-v39}
 \begin{gathered}
 \kappa^2=-tQ/\beta,\qquad
 p\kappa=-\kappa p,\qquad r\kappa=-\kappa r,\\
 p^2=r^2=pr=rp=0
 \end{gathered}
\end{equation}
である。従って $A_\Sigma$ は $R_\Sigma$ 上自由で階数 $6d^2$、
その中心は $R_\Sigma$ であり、
\[
 \operatorname{rad}(A_\Sigma)^2=0,
 \qquad
 A_\Sigma/\operatorname{rad}(A_\Sigma)\simeq M_d(L_\Sigma).
\]
ここで根基はJacobson根基と素根基のいずれにも一致する。
\end{theorem}

\begin{proof}
$A_\Sigma$ 内で
\[
 z=-\frac{\beta x^2}{tQ},\qquad
 f_i=\frac1d\sum_{k=0}^{d-1}t^{-2ik}z^k
 \quad(0\le i<d)
\]
と置く。$z^d=1$ であり、$t^2$ は原始 $d$ 乗根なので、
$f_i$ は直交冪等元で総和は $1$ である。
これらは $\kappa$ や $\lambda$ を添加せず、元の $R_\Sigma$ 上で定義される。
$u$ は $f_i$ から $f_{i-1}$ へ、$v$ は逆向きへ移り、添字は $d$ を法として読む。
さらに
\[
 vu f_i=a_i f_i,\qquad a_i=Q(t^{2i}-1).
\]
$a_0=0$ だが $a_i$ は $1\le i<d$ で $R_\Sigma$ の単元である。
従って、この $d-1$ 本の辺は全ての $f_i$ を結ぶ可逆な鎖を与える。
鎖に沿う行列単位を取ると
\[
 A_\Sigma\simeq M_d(f_0A_\Sigma f_0)
\]
となる。

角代数で $\kappa=xf_0$、$p=f_0u^df_0$、$r=f_0v^df_0$ と置く。
$t^d=-1$ と半周ノルム
\[
 u^dv^d=v^du^d=(-1)^dQ^d+\beta^d\Xi=(-1)^dF
\]
から、式\eqref{eq:singular-corner-relations-v39}を得る。
一般化Weyl代数の次数分解で角に残る次数は $d$ の倍数だけであり、
$u^{2d}=v^{2d}=0$ なので、角は
\[
 1,\ \kappa,\ p,\ \kappa p,\ r,\ \kappa r
\]
で $R_\Sigma$ 加群として生成される。

追加の関係がないことを、逆向きの表示によって確かめる。
式\eqref{eq:singular-corner-relations-v39}で定めた六次元自由 $R_\Sigma$ 代数を
$\mathscr C_\Sigma$ とし、$0\le i,j<d$ の行列単位を $E_{ij}$ とする。
\[
 c_d=\prod_{i=1}^{d-1}a_i=(-1)^{d-1}dQ^{d-1}\in R_\Sigma^\times
\]
と置き、$M_d(\mathscr C_\Sigma)$ 内で
\begin{equation}
\label{eq:singular-descent-matrices-v39}
 \begin{aligned}
 x&=\sum_{i=0}^{d-1}t^i\kappa E_{ii},\\
 u&=\sum_{i=1}^{d-1}E_{i-1,i}+pE_{d-1,0},\\
 v&=\sum_{i=1}^{d-1}a_iE_{i,i-1}+c_d^{-1}rE_{0,d-1}
 \end{aligned}
\end{equation}
と定める。これらは $ux=txu$、$vx=t^{-1}xv$、
$vu=-Q-t^{-1}\beta x^2$、$uv=-Q-t\beta x^2$ を満たす。
折り返しの二項では $t^d=-1$ と $p\kappa=-\kappa p$、$r\kappa=-\kappa r$ を使う。
また $u^d=pI_d$、$v^d=rI_d$、$x^{2d}=\Xi I_d$ であるから、
これは $A_\Sigma$ の表現になる。
$f_i$ は $E_{ii}$ に移り、内部の可逆辺から全ての行列単位が得られる。
$f_0$ の角では $\kappa,p,r$ がそのまま回収される。
従って、先ほど得た六元の生成系は自由基底であり、表示した二つの構成は互いに逆である。

平方零イデアルを除く商は $M_d(L_\Sigma)$ である。
$L_\Sigma\simeq\mathbb C[\kappa,\kappa^{-1},J]$ は整域かつJacobson環であり、
そのJacobson根基は零なので、表示した平方零イデアルが両方の根基となる。
中心を求めるには $\mathscr C_\Sigma$ 内で考えればよい。
$\kappa$ と可換である条件は平方零部分の係数を零にし、
残った $L_\Sigma$ の元が $p,r$ と可換である条件は $\iota$ 不変性である。
従って中心は $L_\Sigma^\iota=R_\Sigma$ である。
\end{proof}

この定理は、中心の開部分全体に対する行列環表示を主張していない。
式\eqref{eq:singular-descent-algebra-v39}の対象は、特異部分への制限 $A_\Sigma$ である。
特に、閉点ファイバーの二頂点を、大域的な中心冪等元で分離することはできない。

\subsection{二つの単純加群は連結な被覆をなす}

\begin{corollary}[単純加群の二重被覆と大域的表示族]
\label{cor:simple-cover-v39}
$A_\Sigma$ の単純加群の同型類は、連結な有限エタール二重被覆
\begin{equation}
\label{eq:simple-cover-map-v39}
 \Omega=\operatorname{Spec}L_\Sigma
 \simeq\mathbb G_{m,\kappa}\times\mathbb A^1_J
 \longrightarrow
 \Sigma\simeq\mathbb G_{m,Q}\times\mathbb A^1_J,
 \qquad Q=-\frac{\beta}{t}\kappa^2
\end{equation}
の閉点でパラメータ化される。
各単純加群の次元は $d$ であり、$\Omega$ 上には各閉点で対応する単純加群を与える
階数 $d$ の大域的表示族が存在する。
複素解析的に $Q$ が原点を一周すると、二つの単純加群は交換される。
\end{corollary}

\begin{proof}
単純加群は平方零根基で零化されるため、
$M_d(L_\Sigma)$ の単純加群と同じである。
これが式\eqref{eq:simple-cover-map-v39}のパラメータ化を与える。
$L_\Sigma$ は表示したLaurent多項式環なので被覆は連結であり、
$2\kappa$ が単元なのでエタールである。

この表示族は式\eqref{eq:singular-descent-matrices-v39}で $p=r=0$ と置けば得られる。
内部の全ての $a_i$ が単元なので、各閉点で得られる $d$ 次元表示は単純である。
元の座標環 $S_H$ に対しても、$w=x^{-1}(J-H)$ とすれば同じ表示が得られる。
この表示では
\begin{equation}
\label{eq:simple-trace-coordinate-v39}
 \operatorname{Tr}(x)
 =\kappa\sum_{i=0}^{d-1}t^i
 =\frac{2\kappa}{1-t}.
\end{equation}
従って二枚の被覆は表示のトレースによって区別される。
最後に $\kappa^2=-tQ/\beta$ から、一周の解析的継続は
$\kappa\mapsto-\kappa$ であり、二つの単純加群を交換する。
\end{proof}

前節の $2d$ 次被覆との関係も明示できる。
$\lambda^{2d}=\Xi$ を $\Sigma$ に制限すると、
\begin{equation}
\label{eq:sigma-cover-components-v39}
 \Sigma\times_Z\operatorname{Spec}Z[\lambda]/(\lambda^{2d}-\Xi)
 =\coprod_{i=0}^{d-1}
 V\bigl(Q+t^{2i-1}\beta\lambda^2\bigr).
\end{equation}
各成分は $\kappa=t^i\lambda$ により $\Omega$ と同型である。
被覆変換 $\lambda\mapsto t\lambda$ は $d$ 成分を巡回し、
$d$ 回施すと各成分内で $\lambda\mapsto-\lambda$ となる。
この最後の作用が単純加群を交換するモノドロミーである。
単純加群の被覆を一律に $\Xi$ の平方根と同一視してはならない。
例えば $d$ が偶数なら $\Xi=(-1)^{d+1}(Q/\beta)^d$ は
$R_\Sigma$ 内の平方であるが、式\eqref{eq:simple-cover-map-v39}はなお連結である。

\subsection{Azumaya代数のBrauer類と冪根の位数}

以下の主張は偶数位数だけでなく、任意の有限位数 $\ell\ge3$ に対して成立する。
$W\subset\operatorname{Spec}Z(A)$ をAzumaya部分とする。
奇数位数では $W=\operatorname{Spec}Z(A)$、偶数位数では
$W=\operatorname{Spec}Z(A)\setminus\Sigma$ である。

\begin{theorem}[正確なBrauer周期と最小分裂次数]
\label{thm:brauer-period-v39}
中心の分数体とその上の代数は
\begin{equation}
\label{eq:generic-symbol-v39}
 \begin{gathered}
 K=\mathbb C(Q,J,\Xi,U),\\
 D=A\otimes_ZK
 \simeq
 \frac{K\langle x,u\rangle}
 {(x^\ell-\Xi,\ u^\ell-U,\ ux-txu)}
 \end{gathered}
\end{equation}
である。$D$ は中心 $K$ 上の斜体であり、
\[
 \operatorname{ind}(D)=\operatorname{per}(D)=\ell.
\]
$W$ 上のAzumaya代数 $\mathcal A$ のBrauer類も正確に位数 $\ell$ を持つ。
有限エタール被覆
\begin{equation}
\label{eq:brauer-splitting-cover-v39}
 W'=W\times_Z\operatorname{Spec}Z[\lambda]/(\lambda^\ell-\Xi)
 \longrightarrow W
\end{equation}
は連結で次数 $\ell$ であり、$\mathcal A$ を分裂させる。
すなわち階数 $\ell$ のベクトル束 $\mathcal P$ が存在して
$\mathcal A|_{W'}\simeq\mathcal End_{\mathcal O_{W'}}(\mathcal P)$ となる。
さらに $D$ を分裂させる任意の有限体拡大 $K'/K$ の次数は $\ell$ で割り切れる。
従って式\eqref{eq:brauer-splitting-cover-v39}は、一般点における有限分裂次数として最小である。
\end{theorem}

\begin{proof}
中心の表示から $V$ は $U$ を可逆にすれば消去できるので、分数体は表示した $K$ となる。
この体上では $u$ が可逆であり、
$v=u^{-1}(-Q-t\beta x^2)$ と書ける。
従って $D$ は式\eqref{eq:generic-symbol-v39}の巡回代数である。
$A=B[x^{-1}][J]$ は整域なので、その中心分数環 $D$ も整域である。
中心上の次元は $\ell^2$ だから、$D$ は次数 $\ell$ の斜体である。

周期を確認するため、$E=K(x)$、$x^\ell=\Xi$ と置く。
$T^\ell-\Xi$ は $\Xi$ に関するEisenstein多項式なので、
$E/K$ は次数 $\ell$ の巡回拡大で、$\tau(x)=tx$ が生成元となる。
ここでは巡回代数の規約を $ua=\tau(a)u$、$u^\ell=U$ とする。
巡回代数の分裂条件は
\begin{equation}
\label{eq:cyclic-norm-criterion-v39}
 [(E/K,\tau,b)]=0
 \quad\Longleftrightarrow\quad b\in N_{E/K}(E^\times)
\end{equation}
である\cite[\S1, equation~(2)]{MinacSchultzSwallowDescent39}。
この条件は、分裂表示を $E$ 上一次元と見て、
$u$ の作用を $z\mapsto c\tau(z)$ と書けば直接分かる。
$u^\ell$ は $N_{E/K}(c)$ 倍となり、逆にノルム表示から分裂表現を作れる。
また巡回代数のBrauer類の積はパラメータの積に対応する。
例えば巡回群の標準的な繰上がり $2$-cocycle を掛け合わせれば、
$[D]^r=[(E/K,\tau,U^r)]$ を得る。

$k_0=\mathbb C(Q,J,\Xi)$ と置けば $K=k_0(U)$、$E=k_0(x)(U)$ である。
$U$ に関する位数 $v_U$ は全ての $\tau^i$ で不変なので、任意の $c\in E^\times$ に対して
\[
 v_U\bigl(N_{E/K}(c)\bigr)
 =\sum_{i=0}^{\ell-1}v_U(\tau^i(c))
 =\ell\,v_U(c).
\]
従って $U^r$ がノルムなら $\ell\mid r$ である。
一方 $U^\ell=N_{E/K}(U)$ なので、式\eqref{eq:cyclic-norm-criterion-v39}より
$\operatorname{per}(D)=\ell$ となる。
この付値は有理関数体上で用いており、$W$ 内の因子を指定しているわけではない。

次数 $\ell$ のAzumaya代数のBrauer類は $\ell$ で零化される。
実際、エタール局所的な行列表示の貼り合わせを $\mathrm{GL}_\ell$ へ持ち上げると、
三重共通部分のスカラー $2$-cocycle の $\ell$ 乗は、その行列式の境界となる。
$\mathcal A$ の類の一般点への制限は周期 $\ell$ の $[D]$ なので、
$W$ 上での周期も正確に $\ell$ である。

次に $\lambda^\ell=\Xi$ を添加する。
$\Xi$ が単元なのでこれは有限エタール拡大であり、一般点で $E/K$ を与えるから連結である。
$x=\lambda$ に対応する冪等元を $e_0$ とする。
一般化Weyl代数の次数分解で $e_0$ の角に残る次数は $\ell$ の倍数であり、
その角は中心 $Z[\lambda]/(\lambda^\ell-\Xi)$ と一致する。
$W'$ 上では、巡回quiverの可逆辺が全頂点を結ぶため $e_0$ はfullである。
従って $\mathcal P=\mathcal A|_{W'}e_0$ は階数 $\ell$ のベクトル束で、
表示した自己準同型環の同型を得る。
これはベクトル束の自由性、すなわち大域的な定数行列環表示を必要としない。

最後に、$D\otimes_KK'\simeq M_\ell(K')$ なら $K'^\ell$ は左 $D$ 加群となる。
その $K$ 次元は $\ell[K':K]$ である。
一方、斜体 $D$ 上の有限次元加群の $K$ 次元は $\dim_KD=\ell^2$ の倍数なので、
$\ell\mid[K':K]$ となる。
\end{proof}

\begin{corollary}[完成局所NCCRの大域降下を妨げるもの]
\label{cor:global-End-obstruction-v39}
元の中心 $Z$ 上で、非零の有限反射的加群 $M$ による
\[
 A\simeq\operatorname{End}_Z(M)
\]
という表示は存在しない。
さらに $A$ は、この形の自己準同型環と $Z$ 上でMorita同値にもならない。
従って前節の完成局所的な非可換クレパント解消は、
元の中心上での自己準同型環としての非可換クレパント解消を意味しない。
\end{corollary}

\begin{proof}
分数体 $K$ へ移れば
$\operatorname{End}_Z(M)\otimes_ZK=\operatorname{End}_K(M\otimes_ZK)$ は分裂した行列環である。
これは非零のBrauer類 $[D]$ と矛盾する。
Morita同値の場合も一般点のBrauer類を保存するため、同じ矛盾を得る。
\end{proof}

特異部分上の単純加群の二重被覆と、滑らかな部分のBrauer類は異なる情報を記録している。
前者の次数は全ての偶数位数で二であり、平方零角代数の型も共通する。
後者の正確な周期は $\ell$ であり、冪根の位数を回収する。
本節の主張は中心上の代数とその加群に関するものであり、
次数付き座標環や $\operatorname{qgr}$ 圏、導来圏の同値分類へは追加の議論なしに移していない。

\section{境界のアフィン改変と混合項の復元}
\label{sec:endpoint-boundary-v39}

前節までの $\Xi\ne0$ の中心開部分では、混合項 $H$ が消えていた。
本節では別の中心元 $U=u^\ell$ を可逆にする。
この開部分は $\Xi=0$ の境界を含み、混合項は一つの有限巡回ノルムに現れる。
そのノルムは中心のアフィン改変、Azumaya部分、および全ての単純加群の次元を支配する。
一般化Weyl代数の有限軌道を切断点で調べる方法は
Drozd--Guzner--Ovsienko\cite[\S\S2,4]{DGOBoundary39}に遡る。
ここではその方法を本族の中心局所化へ適用し、中心ファイバー全体と混合項の復元まで記述する。
一の冪根における中心とAzumaya部分の明示計算の別の例として
Cooney--Ganev--Jordan\cite{CGJBoundary39}も参照されたい。

\subsection{別の中心開部分における一般化Weyl表示}

$\ell\ge3$、$t$ を原始 $\ell$ 乗根、$\beta\ne0$ とする。
三変数部分と中心元の記号は前節と同じであり、
\[
 ux=txu,\qquad vx=t^{-1}xv,\qquad
 uv=-Q-t\beta x^2
\]
を満たす。$S_H=B[w;\sigma,\delta_H]$、$J=xw+H$ と置く。
本節の代数的表示は任意の重み $m\ge1$ と $H\in B_{m+1}$ に成立する。
端点の全変形への適用では $m=2j\ge4$ とする。

$U=u^\ell$ は $B$ の中心にあり、$\sigma(U)=U$、$\delta_H(U)=0$ なので
$S_H$ の中心にも属する。$C_H=S_H[U^{-1}]$ と書く。
$u$ が可逆になるため
\[
 v=u^{-1}(-Q-t\beta x^2),\qquad
 B[U^{-1}]=\mathbb C[Q,u,u^{-1}][x;\tau],\quad \tau(u)=t^{-1}u.
\]
従って一意に
\begin{equation}\label{eq:boundary-H-gauge-v39}
 H=h(Q,u)+xG,\qquad
 h\in\mathbb C[Q,u,u^{-1}],\quad G\in B[U^{-1}]
\end{equation}
と書ける。$h$ は $H$ を $x=0$、$v=-Q/u$ で評価したものである。
$W=w+G$ と置くと $J=xW+h$ である。

\begin{theorem}[境界を含む一般化Weyl代数]
\label{thm:boundary-GWA-v39}
$R=\mathbb C[Q,J,u,u^{-1}]$、$\tau(u)=t^{-1}u$、$\tau(Q)=Q$、$\tau(J)=J$ と置く。
$C_H$ は $R$ 上の一般化Weyl代数で、生成元 $x,W$ と関係
\begin{equation}\label{eq:boundary-GWA-v39}
 xr=\tau(r)x,\quad Wr=\tau^{-1}(r)W,\quad
 xW=J-h(Q,u),\quad Wx=J-h(Q,tu)
\end{equation}
によって表示される。
\end{theorem}

\begin{proof}
$H=h+xG$ をOre導分に代入すると、
\[
 \delta_H(b)=\delta_h(b)+\sigma(b)G-Gb.
\]
従って $W=w+G$ による変数変換後は導分が $\delta_h$ となる。
$h$ は $Q,u^{\pm1}$ と可換なので、$WQ=QW$、$Wu=tuW$ である。
さらに $Wx=xW+h-h(Q,tu)$ から最後の二関係を得る。
$J$ のOre次数に関する最高係数は $x$ の冪なので、$R$ は $C_H$ に単射に入る。
一般化Weyl代数の正規形
\[
 R\ \oplus\,\bigoplus_{k\ge1}Rx^k
 \ \oplus\,\bigoplus_{k\ge1}RW^k
\]
は $\mathbb C[Q,J,u^{\pm1}][x,x^{-1};\tau]$ 内で一次独立である。
実際 $W=x^{-1}(J-h)$ と書け、各 $x$ 次数の係数は零でない。
これが関係式の完全性を与える。
\end{proof}

\subsection{巡回ノルムによる全中心}

$\Xi=x^\ell$、$\Omega=W^\ell$ とし、
\begin{equation}\label{eq:boundary-spectral-norm-v39}
 P_h(Q,J,U)=\prod_{i=0}^{\ell-1}\bigl(J-h(Q,t^iu)\bigr),\qquad u^\ell=U
\end{equation}
と置く。これは $\mathbb C[Q,U,U^{-1}][J]$ のモニックな $\ell$ 次多項式である。
同じ式は、$\mathbb C[Q,U^{\pm1}]$ 上自由な階数 $\ell$ の環
$\mathbb C[Q,u^{\pm1}]$ における乗法行列 $M_h$ を使って
$P_h=\det(JI-M_h)$ と計算できる。

\begin{theorem}[中心とアフィン改変]
\label{thm:boundary-center-v39}
次の等式が成立する。
\begin{equation}\label{eq:boundary-center-v39}
 Z(C_H)=Z(S_H)[U^{-1}]
 =\frac{\mathbb C[Q,J,U,U^{-1},\Xi,\Omega]}{(\Xi\Omega-P_h)}.
\end{equation}
$C_H$ はこの中心上有限で、PI次数は $\ell$ である。
$D=\mathbb C[Q,J,U^{\pm1},\Xi]$ と置くと、この中心は
\[
 D[P_h/\Xi]\subset D[\Xi^{-1}]
\]
であり、イデアル $(\Xi,P_h)$ と因子 $\Xi=0$ によるアフィン改変である。
\end{theorem}

\begin{proof}
一般化Weyl次数を $\deg x=1$、$\deg W=-1$ とする。
$u$ と可換な成分の次数は $\ell$ の倍数であり、その係数が $x$ と可換である条件は
$\tau$ 不変性である。よって中心は
$R^\tau=\mathbb C[Q,J,U^{\pm1}]$ と $x^\ell,W^\ell$ で生成される。
式\eqref{eq:boundary-GWA-v39}から $\Xi\Omega=P_h$ を得る。
正規形の直和により他の関係はない。
$u$ の指数を $\ell$ で割った余りと $x,W$ の $\ell$ 未満の冪が中心上の有限生成系を与える。
中心の分数体上では $x,u$ が可逆で、$x^iu^j$、$0\le i,j<\ell$ が基底である。
従って中心上の次元は $\ell^2$ である。

中心局所化後の元 $bU^{-n}$ が全ての生成元と可換なら、$U$ の正則性から
$b$ 自身がそれらと可換である。これが $Z(C_H)=Z(S_H)[U^{-1}]$ を与える。
最後に $P_h$ は $\Xi$ に依存しない零でない多項式なので、
$\Xi\Omega-P_h$ は既約であり、$\Omega\mapsto P_h/\Xi$ による表示は単射である。
\end{proof}

$\Xi U\ne0$ の重なりでは、前節の中心と
\[
 \Omega=P_h/\Xi,\qquad
 V=\frac{\prod_{i=0}^{\ell-1}(-Q-t^{2i+1}\beta x^2)}{U}
\]
で貼り合わさる。従って混合項が変えるのは、同じ関数体と同じ $\Xi\ne0$ の開部分から
境界へ延ばす際の改変中心である。$u,v$ を交換すれば $V\ne0$ における同様の表示を得る。

\subsection{全中心ファイバーとAzumaya部分}

中心の閉点 $p=(Q_0,J_0,U_0,\Xi_0,\Omega_0)$ を取る。
$U_0\ne0$ なので $\lambda^\ell=U_0$ を選べる。
巡回軌道上の切断点を
\begin{equation}\label{eq:boundary-cuts-v39}
 \mathcal B_p=\{i\in\mathbb Z/\ell\mathbb Z:
 J_0=h(Q_0,t^i\lambda)\},\qquad r=|\mathcal B_p|
\end{equation}
とする。$\Xi_0=\Omega_0=0$ なら $r\ge1$ である。
この $r$ 個の辺で巡回グラフを切って得る鎖の頂点数を
$d_1,\ldots,d_r$ と書く。従って $d_i\ge1$、$\sum_i d_i=\ell$ である。

$r\ge2$ に対し、$\mathcal R_r$ を次の有限次元代数とする。
$r$ 頂点を巡回的に並べ、各隣接対に両向きの矢を置く。
逆向きへの折返しを零とし、各向きに一周する長さ $r$ の経路も零とする。
$r=2$ では各向きに二本の別の矢があると解釈する。
$\mathcal R_1=\mathbb C$ と置く。
その基底は頂点と、各向きに進む長さ $1,\ldots,r-1$ の経路であり、
$\dim\mathcal R_r=r(2r-1)$ である。

\begin{theorem}[切断間隔による全ファイバー分類]
\label{thm:boundary-fibers-v39}
$\Xi_0\ne0$ または $\Omega_0\ne0$ の場合、中心ファイバーは $M_\ell(\mathbb C)$ である。
$\Xi_0=\Omega_0=0$ の場合、中心ファイバーは
\begin{equation}\label{eq:boundary-inflation-v39}
 \operatorname{End}_{\mathcal R_r}
 \left(\bigoplus_{i=1}^r(e_i\mathcal R_r)^{\oplus d_i}\right)
\end{equation}
と同型である。ここで $e_i\mathcal R_r$ は右加群である。その単純加群はちょうど $r$ 個で、次元は $d_1,\ldots,d_r$ である。
全次元と根基の次元は
\[
 \dim(C_H\otimes\kappa(p))=2\ell^2-\sum_i d_i^2,\qquad
 \dim\operatorname{rad}(C_H\otimes\kappa(p))
 =2\left(\ell^2-\sum_i d_i^2\right).
\]
根基の冪零指数は正確に $r$ である。$r=1$ では根基は零と解釈する。

Azumaya部分は中心の滑らかな部分と一致する。
非Azumaya部分は、正確に $\Xi=\Omega=0$ かつ $r\ge2$ の点である。
\end{theorem}

\begin{proof}
有限エタール拡大 $\lambda^\ell=U$ の後、$u$ は $\ell$ 個の相異なる固有冪等元に分解する。
$x,W$ は巡回グラフの互いに逆向きの矢となり、各辺の往復積は
$J-h(Q,t^i\lambda)$ である。$\Xi$ または $\Omega$ が可逆なら、片向きの全ての矢が可逆なので
全頂点が同型になり、局所的な行列環表示を得る。

両者が零の閉ファイバーでは、往復積が非零の辺を可逆な同型として縮約する。
残る $r$ 頂点が鎖の成分を表し、その重複度は $d_i$ である。
切断辺の往復積は零であり、$x^\ell=W^\ell=0$ が両向きの一周を零にする。
向きを変える経路は折返しを含むため零となり、残る経路は表示した基底に限られる。
これが式\eqref{eq:boundary-inflation-v39}への次元の上界を与える。

逆に、各鎖の内部に $d_i\times d_i$ 行列単位を置き、切断辺を $\mathcal R_r$ の矢で置き換える。
元の非零往復積は、逆向きの内部矢をその非零スカラーで掛ければ実現できる。
全ての関係が成立し、生成される像は行列単位と全ての許された経路を含む。
従って表示した代数への全射があり、上界と下界が一致する。
異なる頂点間の経路は各向きに一本ずつあり、同じ頂点には恒等経路だけがあるため
次元公式を得る。正の経路長が根基であり、最大長は $r-1$ なので冪零指数も従う。

最後に、$r=1$ なら残る $\ell-1$ 本の辺は近傍でも可逆であり、
その縮約によって中心上の行列環表示が得られる。これは中心上の局所自由性も示す。
$r\ge2$ ならファイバーは二つ以上の単純加群を持つのでAzumayaではない。
中心の滑らかさはエタール基底変換後に判定できる。
方程式 $\Xi\Omega=\prod_i(J-h_i)$ において、$\Xi$ または $\Omega$ が非零なら滑らかである。
両者が零で $r=1$ なら $J$ 微分が非零、$r\ge2$ なら全ての一階偏微分が零となる。
よって両部分は一致する。
\end{proof}

\subsection{標識付き境界からの混合項の復元}

\begin{proposition}[巡回ノルムの復元性]
\label{prop:boundary-recovery-v39}
$Q,U,J$ の標識を固定する。二つの Laurent 多項式 $h,h'$ について
\[
 P_h=P_{h'}\quad\Longleftrightarrow\quad
 h'(Q,u)=h(Q,t^ku)\text{ for some }k\in\mathbb Z/\ell\mathbb Z.
\]
従って、モニックな境界方程式 $P_h=0$ は $H\bmod x$ を巡回作用の軌道まで決定する。
\end{proposition}

\begin{proof}
$E=\mathbb C(Q,u)$、$K=\mathbb C(Q,U)$、$U=u^\ell$ は巡回Galois拡大である。
$P_h=P_{h'}$ なら $E$ 内で
\[
 0=P_h(h')=\prod_i\bigl(h'-h(Q,t^iu)\bigr).
\]
$E$ は体なので一つの因子が零である。逆は積の因子の置換から従う。
\end{proof}

この命題は、基底変数と $J$ を固定した境界方程式についての主張である。
標識を忘れた中心環や全射影圏の同型分類までは主張していない。
中心元を $H$ に加えると $J$ の平行移動が生じるので、この標識の指定は必要である。

さらに $G_h\subset\langle t\rangle$ を $h$ の安定化群、$g=|G_h|$ とすると、
$P_h=P_{\mathrm{min}}^g$ であり、$P_{\mathrm{min}}$ は $K$ 上の最小多項式である。
既約な被約境界因子の正規化は
\[
 \operatorname{Spec}\mathbb C[Q,u^{\pm1}]^{G_h}
 =\operatorname{Spec}\mathbb C[Q,u^{\pm g}]
\]
で与えられる。実際、その環は $\mathbb C[Q,U^{\pm1},h]$ 上有限、正規であり、
分数体はGalois理論から $K(h)=E^{G_h}$ に一致する。
一般の境界点では切断点数が $g$、各単純加群の次元が $\ell/g$ となる。
安定化群以上の偶然の値の一致が起きる部分で、さらに単純加群が分裂する。

\begin{theorem}[一般の境界における巡回降下と単純加群の被覆]
\label{thm:boundary-cyclic-descent-v39}
$g=|G_h|$、$n=\ell/g$、$z=u^g$ とし、被約境界因子の正規化上で
\[
 h_i(z)=h(Q,t^iu)\quad(0\le i<n),\qquad
 \Delta=\prod_{i=1}^{n-1}(h_0-h_i),\qquad
 A=\mathbb C[Q,z,z^{-1},\Delta^{-1}]
\]
と置く。$n=1$ の積は $1$ とする。
$U=z^n$、$J=h_0$、$\Xi=\Omega=0$ により中心から $A$ への写像を定める。
$L=A[\kappa]/(\kappa^g-z)$ とし、$\sigma(\kappa)=t^{-n}\kappa$ と置く。
次の代数を定める。
\[
 \mathcal T_g=
 \frac{L\langle P,R\rangle}
 {(Pa-\sigma(a)P,\ Ra-\sigma^{-1}(a)R,\ PR,\ RP,\ P^g,\ R^g:\ a\in L)}.
\]
$g=1$ では $\mathcal T_1=A$ と解釈する。このとき
\[
 C_H\otimes_{Z(C_H)}A\simeq M_n(\mathcal T_g).
\]
$\mathcal T_g$ の正の $P,R$ 次数の部分は冪零イデアルであり、商は $L$ である。
従って、この部分上の単純加群の同型類は連結な有限エタール巡回被覆
\[
 \operatorname{Spec}L\longrightarrow\operatorname{Spec}A
\]
で標識される。被覆の次数は $g$、単純加群の次元は $n$ であり、
商写像 $M_n(\mathcal T_g)\to M_n(L)$ がその族の実際の表現を与える。
\end{theorem}

\begin{proof}
$h$ の各 $u$ 指数は $g$ の倍数なので、$h_i$ は $z$ のLaurent多項式である。
$\Delta\ne0$ は安定化群以上の値の一致を除く条件である。
$u^g$ の固有値 $t^{gi}z$ に対応する $n$ 個の冪等元を取る。
$x$ は添字を一つ進め、$W$ は逆に進める。
添字 $0$ 以外の往復積 $b_i=h_0-h_i$ は $A$ の単元なので、
これら $n$ 個の冪等元は同型である。
添字零の角では $u=\kappa$、$P=x^n$、$R=W^n$ と置ける。
正規形のうち角に残る次数は $n\mathbb Z$ であり、
$P^g=R^g=0$、$PR=RP=0$、および表示した係数との交換則が成立する。

追加関係がないことを明示的にも確認できる。
$M_n(\mathcal T_g)$ 内で
\[
 u=\operatorname{diag}(\kappa,t\kappa,\ldots,t^{n-1}\kappa),\qquad
 x=\sum_{i=0}^{n-2}E_{i+1,i}+P E_{0,n-1},
\]
\[
 W=\sum_{i=1}^{n-1}b_iE_{i-1,i}
       +\Delta^{-1}R E_{n-1,0}
\]
と置けば全関係が成立し、内部行列単位と角の全生成元を得る。
$n=1$ では $x=P$、$W=R$ と読む。
逆向きの角写像と合わせて同型が従う。

$L$ は $A$ 上有限エタールである。分数体上で $\kappa^g-z$ は
$z$ の付値によって既約なので、被覆は連結である。
冪零イデアルを割った環は $M_n(L)$ であり、各閉点では $g$ 個の
$n$ 次元単純加群になる。$\kappa\mapsto t^n\kappa$ がそれらを巡回的に交換する。
\end{proof}

この定理は、被約境界因子の正規化から偶然の追加衝突を除いた部分に関する表示である。
衝突点を含む全中心ファイバーについては、定理\ref{thm:boundary-fibers-v39}の
任意の切断間隔による分類を使う。

\subsection{同じ関数斜体と開部分を持つ異なる境界}

$m=4$、$\ell=5$ とする。二つの許された混合項
\[
 H_0=u^5,\qquad H_1=u^4v
\]
を比較する。$h_0=u^5=U$、$h_1=-Qu^3$ なので、中心はそれぞれ
\begin{equation}\label{eq:boundary-order5-example-v39}
 Z(C_{H_0})=\frac{\mathbb C[Q,J,U^{\pm1},\Xi,\Omega]}
 {(\Xi\Omega-(J-U)^5)},\qquad
 Z(C_{H_1})=\frac{\mathbb C[Q,J,U^{\pm1},\Xi,\Omega]}
 {(\Xi\Omega-J^5-Q^5U^3)}.
\end{equation}
前者の特異部分は $\Xi=\Omega=J-U=0$ で次元二、後者は
$\Xi=\Omega=J=Q=0$ で次元一である。従ってこの二つの中心開部分は環として同型でない。
前者の一般の境界点には五つの一次元単純加群があり、
後者の $Q\ne0$ の境界点では一つの五次元単純加群がある。
一方、両者の関数斜体と $\Xi\ne0$ の中心局所化は前節の定理により同型である。

より一般に $h=cQ^a u^n$、$cQ^a\ne0$、$g=\gcd(\ell,n)$ とすると
\[
 P_h=\left(J^{\ell/g}-c^{\ell/g}Q^{a\ell/g}U^{n/g}\right)^g.
\]
境界の切断点は等間隔であり、単純加群は $g$ 個、各次元は $\ell/g$ となる。
上の例とこの一般式を、巡回ノルムと乗法行列の行列式という独立な二方法で検証した。

本節により、混合項の情報が境界への延長に現れるという前版の見通しを、
二つの中心開部分の間の具体的な貼合せと復元定理として実現できた。
$\Xi,U,V$ が全て零となる残りの閉部分、および標識を忘れた大域的分類は今後の問題である。

\section{全中心、最深点の単純加群と完備局所Brauer障害}
\label{sec:global-center-v40}

本節では、混合項 $H=f(Q,u)+xG$ の部分族について全中心を決定する。
そのうち $H=0$ と同型な部分族では、全ての閉中心ファイバーも分類できる。
最深点では、単純加群が一つに合流する一方、完備局所分数体上のBrauer指数は
最大値 $\ell$ を保つ。この点は、$\Xi\ne0$ の閉点完備化で得られた
分裂した局所模型と異なる。

\subsection{対象と全中心の完全交叉表示}

引き続き $t$ は原始 $\ell$ 乗根、$\ell\ge3$、$\beta\ne0$ とする。
$H=0$ の座標環を
\[
 S_0=B[w;\sigma],\qquad
 B=\frac{\mathbb C\langle x,u,v\rangle}
 {(ux-txu,\ vx-t^{-1}xv,\ uv-vu-(t^{-1}-t)\beta x^2)}
\]
と書く。ここで
\[
 wx=xw,\qquad wu=tuw,\qquad wv=t^{-1}vw,
 \qquad Q=-uv-t\beta x^2,
 \qquad J=xw.
\]
$w$ の次数 $m$ は任意の正整数でよい。
本節の中心と表現については、次数を忘れた代数として扱う。

より一般に、斉次元 $H$ が
\begin{equation}\label{eq:global-gauge-class-v40}
 H=H_c+xG,\qquad H_c\in Z(B)
\end{equation}
と書ける場合も全て以下に含まれる。
実際、$W=w+G$ と置くと $xW+H_c=xw+H$ となり、
$Wb=\sigma(b)W$ が成立する。
$H_c$ の中心性から $\sigma(H_c)=H_c$ であることも従う。
従って $S_H\simeq S_0$ であり、中心の座標 $J$ だけを $J-H_c$ に置き換えればよい。
任意の混合項が式\eqref{eq:global-gauge-class-v40}を満たすとは仮定しない。

\begin{theorem}[全中心と正規性]\label{thm:global-center-v40}
$\Xi=x^\ell$、$\Omega=w^\ell$、$U=u^\ell$、$V=v^\ell$ と置く。
\[
 P_\ell(Q,\Xi)=
 \begin{cases}
  -Q^\ell-\beta^\ell\Xi^2,&\ell\text{ が奇数},\\
  (Q^d+(-1)^d\beta^d\Xi)^2,&\ell=2d
 \end{cases}
\]
とすれば、全中心は
\begin{equation}\label{eq:global-center-v40}
 Z_0:=Z(S_0)=
 \frac{\mathbb C[Q,J,\Xi,\Omega,U,V]}
 {(\Xi\Omega-J^\ell,\ UV-P_\ell(Q,\Xi))}
\end{equation}
である。$Z_0$ は四次元の正規Gorenstein整域であり、二つの方程式による完全交叉である。
$S_0$ は $Z_0$ 上有限で、一般階数は $\ell^2$、PI次数は $\ell$ である。
\end{theorem}

\begin{proof}
$R=\mathbb C[Q,x,w]$、
\[
 \tau(x)=tx,\quad\tau(w)=t^{-1}w,\quad\tau(Q)=Q,
 \qquad a=-Q-t^{-1}\beta x^2
\]
と置く。$S_0$ は $R$ 上の一般化Weyl代数であり、
\[
 ur=\tau(r)u,\quad vr=\tau^{-1}(r)v,\quad vu=a,\quad uv=\tau(a)
\]
と表示される。$R$ が多項式環として入ることは、元のOre基底と
$Q=-uv-t\beta x^2$ の最高 $u,v$ 次項から分かる。
一般化Weyl代数の正規形
\[
 R\ \oplus\bigoplus_{k\ge1}Ru^k\ \oplus\bigoplus_{k\ge1}Rv^k
\]
は $R[u,u^{-1};\tau]$ 内で一次独立である。
これにより関係式の完全性も分かる。

中心元の一般化Weyl次数は $\ell$ の倍数である。
その係数は $\tau$ 不変でなければならないから、中心は
$R^\tau$ と $u^\ell,v^\ell$ で生成される。
$x^aw^b$ が不変である条件は $a-b\equiv0\pmod\ell$ であり、従って
\[
 R^\tau=\mathbb C[Q,\Xi,\Omega,J]/(\Xi\Omega-J^\ell).
\]
$UV=\prod_{i=0}^{\ell-1}(-Q-t^{2i+1}\beta x^2)=P_\ell$ は既に計算した巡回ノルムである。
上の正規形により、この二つ以外の中心関係はない。
$R$ は $R^\tau$ 上有限であり、$u,v$ の指数は $\ell$ 未満に還元できるので、
中心上の有限性も従う。一般階数とPI次数は $\Xi U\ne0$ での基底
$x^iu^j$、$0\le i,j<\ell$ から得る。

表示は整域であり、次元は四である。
二方程式は正則列なのでCohen--MacaulayかつGorensteinである。
直後のJacobian計算で特異部分の次元は高々二と分かる。
従って余次元一で正則であり、Serreの判定により正規である。
\end{proof}

\begin{theorem}[一変数混合項の全中心]
\label{thm:global-polynomial-center-v40}
$f(Q,u)\in\mathbb C[Q,u]$ とし、
$\delta_f(b)=x^{-1}[b,f]$、$S_f=B[w;\sigma,\delta_f]$ を取る。
$J=xw+f(Q,u)$、$\Omega=w^\ell$ と置く。
巡回ノルム
\[
 N_f(Q,U,J)=\prod_{i=0}^{\ell-1}\bigl(J-f(Q,t^iu)\bigr),\qquad U=u^\ell
\]
は $\mathbb C[Q,U][J]$ のモニックな $\ell$ 次多項式である。
全中心は
\begin{equation}\label{eq:global-polynomial-center-v40}
 Z(S_f)=
 \frac{\mathbb C[Q,J,\Xi,\Omega,U,V]}
 {(UV-P_\ell(Q,\Xi),\ \Xi\Omega-N_f(Q,U,J))}.
\end{equation}
これは四次元の正規Gorenstein完全交叉であり、$S_f$ はその上有限でPI次数 $\ell$ である。
同じ結論は $H=f(Q,u)+xG$ に対し、$W=w+G$、$\Omega=W^\ell$ と置いて成立する。
次数付き座標環として用いる場合は $f$ を元の重みで次数 $m+1\ge2$ の斉次式に取る。
\end{theorem}

\begin{proof}
$B/(x)$ は可換なので $\delta_f$ は $B$ に値を持つ $\sigma$ 導分である。
$B$ は整域、$\sigma$ は自己同型であり、従ってOre拡大 $S_f$ も整域である。
$U$ を可逆にした第\ref{sec:endpoint-boundary-v39}節の表示では、
$h=f(Q,u)$、$G=0$ とできる。従って $w^\ell$ はその局所化の中心に属する。
元の $S_f$ は整域であるから、$[w^\ell,b]U^{-k}=0$ なら
$[w^\ell,b]=0$ である。よって $\Omega=w^\ell$ は全代数で中心的である。
他の三つの $\ell$ 乗も中心的であるから、Ore基底の四つの指数を
$\ell$ 未満に還元できる。従って、表示した中心元が生成する部分環上で $S_f$ は有限である。

式\eqref{eq:global-polynomial-center-v40}の右辺を $A_f$ と書く。
まず $A_f$ が整域であることを確認する。
\[
 R_B=\mathbb C[Q,\Xi,U,V]/(UV-P_\ell(Q,\Xi))
\]
は整域である。$N_f$ は $J$ に関してモニックなので、$A_f$ は
$R_B[\Omega]$ 上、$1,J,\ldots,J^{\ell-1}$ を基底とする自由加群である。
特に $\Xi$ は $A_f$ 上の非零因子である。
一方、$\Xi$ を可逆にすれば $\Omega=N_f/\Xi$ と消去でき、
$A_f[\Xi^{-1}]=R_B[\Xi^{-1}][J]$ は整域となる。
従って $A_f$ 自身も整域であり、$\Xi$ の局所化を通して $Z(S_f)$ へ単射に入る。

二方程式は正則列であるから $A_f$ は四次元のGorenstein完全交叉である。
$D(\Xi)$ は定理\ref{thm:root-center-v38}の正規な中心開部分である。
$D(U)$ は定理\ref{thm:boundary-center-v39}のノルム超曲面
$\Xi\Omega=N_f$ となる。
その特異点では $\Xi=\Omega=0$ が必要であり、
$N_f=0$ は $J$ に関してモニックなので特異部分の次元は高々二である。
従ってこの四次元超曲面も正規である。
残る $V(\Xi,U)$ の被約台では
\[
 Q=0,\qquad J=f(0,0)
\]
であり、自由な変数は $\Omega,V$ だけなので次元は二である。
よって $A_f$ は余次元一で正則、かつCohen--Macaulayであり、全体が正規である。

$Z(S_f)$ は有限 $A_f$ 加群 $S_f$ の部分加群なので、$A_f$ 上有限である。
$D(\Xi)$ で両環は一致するから分数体も同じである。
$A_f$ の整閉性により $Z(S_f)=A_f$ となる。
一般点の巡回代数は以前と同じなのでPI次数は $\ell$ である。
最後の変数変換は $W=w+G$ をOre関係へ代入すれば得られる。
\end{proof}

例えば $f=u^k$、$g=\gcd(\ell,k)$ の場合、第二の方程式は
\[
 \Xi\Omega=(J^{\ell/g}-U^{k/g})^g
\]
となる。$f=Q^a u^k$ なら同じ右辺で $U^{k/g}$ を
$Q^{a\ell/g}U^{k/g}$ に置き換える。
この定理は非中心的な混合項も含むが、任意の二変数混合項の全中心まで決定したものではない。
以下の全ファイバー分類は $S_0$ に対して述べる。
最深点の完備局所定理は、その後で $f\in Q\mathbb C[Q,u]$ へ拡張する。

\subsection{全特異部分と全ての中心ファイバー}

$\operatorname{Spec}Z_0$ の原点を $o$ とする。次の閉部分を被約構造で定める。
\begin{equation}\label{eq:global-strata-v40}
 T=V(\Xi,\Omega,J),\qquad
 D=V(Q,\Xi,J,U,V).
\end{equation}
$T$ は $UV=(-Q)^\ell$ を満たす二次元の $A_{\ell-1}$ 型曲面であり、
$D\simeq\mathbb A^1_\Omega$、$T\cap D=\{o\}$ である。
$\ell=2d$ のときは、さらに
\[
 F=Q^d+(-1)^d\beta^d\Xi,\qquad E=V(U,V,F)
\]
と置く。$D\subset E$、$E\cap T=\{o\}$ である。

\begin{theorem}[全Azumaya部分とファイバーの分類]
\label{thm:global-fibers-v40}
$S_0$ のAzumaya部分は $Z_0$ の滑らかな部分と一致する。
被約特異部分は
\[
 (\operatorname{Sing}Z_0)_{\mathrm{red}}=
 \begin{cases}
 T\cup D,&\ell\text{ が奇数},\\
 T\cup E,&\ell=2d.
 \end{cases}
\]
$\mathcal R_r$ は定理\ref{thm:boundary-fibers-v39}の巡回切断代数とする。
閉中心ファイバーは次のように分類される。
\begin{center}
\small
\begin{tabular}{c|c|c|c}
中心の位置&ファイバー&単純加群の次元&根基の冪零指数\\\hline
滑らかな部分&$M_\ell(\mathbb C)$&$\ell$ が一つ&$1$\\
$T\setminus\{o\}$ または $D\setminus\{o\}$&$\mathcal R_\ell$&$1$ が $\ell$ 個&$\ell$\\
$E\setminus D$、$\ell=2d$&$M_d(\mathcal R_2)$&$d$ が二つ&$2$\\
$o$&$\mathcal F_\ell$&$1$ が一つ&$2\ell-1$
\end{tabular}
\end{center}
最後の $\mathcal F_\ell$ の基底と次元は次の定理で与える。
零根基の冪零指数は $1$ としている。
\end{theorem}

\begin{proof}
中心の二方程式のJacobianを用いる。
第一行が零となるのは正確に $\Xi=\Omega=J=0$、すなわち $T$ である。
$\ell$ が奇数のとき、第二行が零となるのは
$Q=\Xi=U=V=0$ であり、中心の方程式から $J=0$ も従う。
これは $D$ である。$\ell=2d$ のとき、$\partial F/\partial\Xi\ne0$ なので
第二行が零となるのは正確に $U=V=F=0$、すなわち $E$ である。
両行の共通する列は $\Xi$ 列だけである。
$\Xi\ne0$ では第一行の $\Omega$ 成分が非零なので、二行が従属になるには第二行が零でなければならない。
$\Xi=0$ では $J=0$ であり、$\Omega\ne0$ の場合も
$Q,U,V$ の各成分を調べると上で列挙した部分以外の従属は生じない。
これで被約特異部分が得られる。

$\Xi\ne0$ のファイバーは定理\ref{thm:root-Azumaya-v38}で既に決定されている。
$\Omega\ne0$ では $w$ が可逆なので $x=Jw^{-1}$ と書ける。
一般化Weyl表示の係数環は $\mathbb C[Q,J,w^{\pm1}]$ となり、
$w$ の巡回軌道は常に長さ $\ell$ である。
$\Xi=0$ の閉点では $J=0$、従って $x=0$ である。
$Q\ne0$ なら往復積 $a=-Q$ が全て非零なので、ファイバーは $M_\ell(\mathbb C)$ になる。
$Q=0$ なら全ての辺が切断され、$U=V=0$ のときの全ファイバーは
$\mathcal R_\ell$ になる。$U$ または $V$ が非零なら一方向が可逆となるので再び全行列環である。

残る $\Xi=\Omega=J=0$ では、原点以外なら $U$ または $V$ が非零である。
例えば $U\ne0$ では $u$ が可逆で、$v=u^{-1}(-Q-t\beta x^2)$ と消去できる。
閉ファイバーでは、$u$ の $\ell$ 個の固有冪等元の間を $x,w$ が逆向きに結び、
$xw=wx=0$、$x^\ell=w^\ell=0$ が成立する。
これは正確に $\mathcal R_\ell$ の表示である。
$V\ne0$ の場合も同様である。
これで原点以外の全ファイバーと非Azumaya性が分かる。
滑らかな部分では、上記の局所化における可逆辺の縮約によって
中心上エタール局所的な行列環表示が得られる。
従って、ファイバーの単純性だけでなく局所自由性を含むAzumaya性が従う。
原点は次の定理により非Azumayaである。
\end{proof}

\begin{proposition}[偶数の場合の特異成分の正規化]
\label{prop:global-singular-normalization-v40}
$\ell=2d$、$\alpha=(-1)^d\beta^d$ とする。
$E$ の正規化は $A_1$ 型曲面
\[
 \widetilde E=\operatorname{Spec}\mathbb C[Q,J,a]/(Qa-J^2)
\]
であり、正規化写像は
\[
 \Xi=-Q^d/\alpha,\qquad \Omega=-\alpha a^d,\qquad U=V=0
\]
で与えられる。
$D\setminus\{o\}$ の各点の逆像は $d$ 点からなる。
\end{proposition}

\begin{proof}
$E$ の方程式は $Q^d\Omega=-\alpha J^{2d}$ である。
$Q\ne0$ で $a=J^2/Q$ と置くと $a^d=-\Omega/\alpha$ であり、
$a$ は元の環上整である。表示した環は同じ分数体を持つ正規な $A_1$ 曲面なので正規化である。
$Q=J=0$、$\Omega\ne0$ では $a^d=-\Omega/\alpha$ が相異なる $d$ 個の解を持つ。
\end{proof}

\subsection{最深点の全代数と単純加群の合流}

原点の中心ファイバーを
\[
 \mathcal F_\ell=S_0/(Q,J,\Xi,\Omega,U,V)
\]
と書く。これを考える際は $x,u,v,w$ の全てに補助次数 $1$ を与える。
この補助次数は元の重み $m$ を置き換えてファイバーの根基を測るためのものである。

\begin{theorem}[最深点の基底と根基]\label{thm:deepest-fiber-v40}
$e_0=1$、$e_k=u^k$ ($k>0$)、$e_k=v^{-k}$ ($k<0$) と置き、
\[
 b_0=\ell,\qquad b_k=\min\{\ell,2(\ell-|k|)\}\quad(0<|k|<\ell)
\]
と定める。次の集合は $\mathcal F_\ell$ の基底である。
\begin{equation}\label{eq:deepest-basis-v40}
 \{x^ae_k:1-\ell\le k\le\ell-1,\ 0\le a<b_k\}
 \ \cup\
 \{w^be_k:1-\ell\le k\le\ell-1,\ 1\le b<\ell\}.
\end{equation}
従って
\begin{equation}\label{eq:deepest-dimension-v40}
 \dim\mathcal F_\ell=
 \begin{cases}
 (7\ell^2-6\ell+1)/2,&\ell\text{ が奇数},\\
 (7\ell^2-6\ell+2)/2,&\ell\text{ が偶数}.
 \end{cases}
\end{equation}
$\mathcal F_\ell$ は局所代数であり、唯一の単純加群は一次元である。
根基は正の補助次数の部分で、その冪零指数は正確に $2\ell-1$ である。
\end{theorem}

\begin{proof}
まず中心関係 $Q=J=\Xi=\Omega=0$ を課し、
\[
 A=\mathbb C[x,w]/(x^\ell,w^\ell,xw),\qquad
 G=A(\tau,-t^{-1}\beta x^2)
\]
と置く。この一般化Weyl代数は
$G=\bigoplus_{k\in\mathbb Z}Ae_k$ という直和分解を持つ。
係数環が整域でなくてもこの正規形は成立する。
実際、$vu=a$、$uv=\tau(a)$ の重なり $uvu$ と $vuv$ は
それぞれ $u a=\tau(a)u$ と $v\tau(a)=av$ により解消され、
係数との交換則も同じ自己同型 $\tau$ によって整合する。

$\tau^\ell=1$ なので、$u^\ell,v^\ell$ は $G$ の中心元である。
$\mathcal F_\ell=G/(u^\ell,v^\ell)$ を一般化Weyl次数ごとに計算する。
$|k|\ge\ell$ の成分は全て消える。
$0\le k<\ell$ では、中心イデアルの第 $k$ 成分は
\[
 u^\ell G_{k-\ell}+v^\ell G_{k+\ell}
 =A x^{2(\ell-k)}e_k+A x^{2\ell}e_k
 =A x^{2(\ell-k)}e_k.
\]
ここで各積の非零定数係数はイデアルを変えず、$x^{2\ell}=0$ を使った。
負の $k$ でも $u,v$ を交換して同様に計算できる。
$k=0$ の成分では追加関係はない。
従って残る第 $k$ 成分は正確に
$A/(x^{2(\ell-|k|)})e_k$ であり、式\eqref{eq:deepest-basis-v40}を得る。

検算に使える左生成元の作用も明記する。
基底の範囲外は零と解釈し、$x,w$ は係数の左乗法で作用させる。
$x^a$ と $w^b$ が混ざる積は零であり、残る作用は
\[
\begin{aligned}
 u(re_k)&=\tau(r)e_{k+1}&& (k\ge0),\\
 u(re_k)&=-t\beta\tau(r)x^2e_{k+1}&& (k<0),\\
 v(re_k)&=\tau^{-1}(r)e_{k-1}&& (k\le0),\\
 v(re_k)&=-t^{-1}\beta\tau^{-1}(r)x^2e_{k-1}&& (k>0).
\end{aligned}
\]
これらは元の表示の関係を満たし、$1$ は循環ベクトルである。
付属コードは元の関係式からこの左正則表示を独立に検証する。

次元は
\[
 (2\ell-1)+2\sum_{k=1}^{\ell-1}
 \bigl(\ell-1+\min\{\ell,2(\ell-k)\}\bigr)
\]
を計算すれば得られる。
全ての基底元の補助次数は高々 $2\ell-2$ であり、
$w^{\ell-1}u^{\ell-1}$ はその最大次数を持つ非零元である。
従って正次数イデアルの冪零指数は $2\ell-1$ であり、
その商は $\mathbb C$ だからこれがJacobson根基である。
\end{proof}

この表示から根基の各層も読み取れる。そのHilbert多項式は
\[
 \sum_{a=0}^{\ell-1}z^a+\sum_{b=1}^{\ell-1}z^b
 +2\sum_{k=1}^{\ell-1}z^k
 \left(\sum_{a=0}^{b_k-1}z^a+\sum_{b=1}^{\ell-1}z^b\right)
\]
である。補助次数の生成元が全て一次なので、係数は
$\dim(\operatorname{rad}^i/\operatorname{rad}^{i+1})$ に一致する。

一次元単純加群そのものも大域的に記述できる。
可換スカラー上では $t\ne\pm1$ より必ず $x=0$ であり、
残る条件は $wu=wv=0$ である。
従って一次元単純加群の被約パラメータ空間は
\[
 \{w=0\}\cup\{u=v=0\}\subset\mathbb A^3_{u,v,w}.
\]
$T\setminus\{o\}$ への中心指標写像は
\[
 (u,v)\longmapsto(Q,U,V)=(-uv,u^\ell,v^\ell)
\]
という連結エタール巡回 $\ell$ 重被覆であり、
$D\setminus\{o\}$ への写像は $w\mapsto\Omega=w^\ell$ という同様の被覆である。
原点ではそれらの単純加群が一つの自明表現に合流する。
これは、中心ファイバーの次元がさらに増大することと両立する。

\subsection{完備局所化後にも消えないBrauer指数}

全中心の原点完備化を $\widehat Z_0$、その分数体を $K_0$ と置く。
次の定理は、完備局所化後の分裂性を一般点での分裂性から推測せず、
新たな付値を構成して判定する。

\begin{theorem}[最深点での完備局所Brauer障害]
\label{thm:deepest-Brauer-v40}
$\widehat Z_0$ は整域である。
\[
 \widehat S_0\otimes_{\widehat Z_0}K_0
 \simeq (\Xi,U)_t
\]
は次数 $\ell$ の巡回斜体であり、そのBrauer周期と指数はいずれも正確に $\ell$ である。
従って $\widehat S_0$ は、$\widehat Z_0$ 上の非零有限反射的加群 $M$ の
$\operatorname{End}_{\widehat Z_0}(M)$ と同型ではなく、
そのような自己準同型環と $\widehat Z_0$ 線形Morita同値にもならない。
\end{theorem}

\begin{proof}
完備冪級数環で次の正の重みを用いる。
\begin{equation}\label{eq:deepest-valuation-weights-v40}
 \operatorname{wt}(Q)=\operatorname{wt}(J)=\operatorname{wt}(\Xi)
 =\operatorname{wt}(U)=\operatorname{wt}(V)=1,
 \qquad \operatorname{wt}(\Omega)=\ell-1.
\end{equation}
$c=(-1)^\ell\beta^\ell$ と置く。
二つの方程式の最低重み項は
\[
 \Xi\Omega-J^\ell,\qquad UV-c\Xi^2
\]
である。これらが付随次数環の全関係を与えることを確認する。
二式を $J^\ell$ と $UV$ を消去する向きに使うと、先頭単項式は互いに素であり、
唯一の重なりは同じ結果に簡約される。
正規単項式は $J$ の指数が $\ell$ 未満で、$U,V$ を同時に含まないものとなる。
第一の簡約は重みを保ち、第二の簡約は重みを下げない。
従って完備化後も各有限重みまで同じ正規形が一意に存在し、
\begin{equation}\label{eq:deepest-associated-graded-v40}
 \operatorname{gr}_{\mathrm{wt}}\widehat Z_0
 =\frac{\mathbb C[Q,J,\Xi,\Omega,U,V]}
 {(\Xi\Omega-J^\ell,\ UV-c\Xi^2)}
\end{equation}
を得る。

右辺は整域である。まず
$A=\mathbb C[Q,J,\Xi,\Omega]/(\Xi\Omega-J^\ell)$ は整域である。
次に $A[U,V]/(UV-c\Xi^2)$ の正規形を $A[U,U^{-1}]$ へ
$V\mapsto c\Xi^2/U$ で写すと、負の $U$ 次数の係数は
非零元 $c\Xi^2$ の冪だけを掛けられるので、単射である。
従って付随次数環は整域であり、完備環も整域である。
最低重みは積について加法的になり、$K_0$ 上の離散付値 $\nu$ を定める。
$\nu(\Xi)=1$ だからその値群は $\mathbb Z$ である。

この付値の剰余体は
\begin{equation}\label{eq:deepest-residue-field-v40}
 \kappa_\nu=\mathbb C(q,j,r),\qquad
 q=\overline{Q/\Xi},\quad j=\overline{J/\Xi},\quad r=\overline{U/\Xi}
\end{equation}
である。実際、付随次数環の分数体で
$V/\Xi=c/r$、$\Omega/\Xi^{\ell-1}=j^\ell$ と書け、
$q,j,r$ は代数的に独立である。

一般点の巡回表示は $\Xi U\ne0$ での表示を体拡大することで得る。
この時点では体拡大後に斜体であるとは仮定しない。
付値完備化 $K_\nu$ を取り、$E=K_\nu(x)$、$x^\ell=\Xi$ と置く。
これはEisenstein多項式による完全分岐巡回拡大で、次数は $\ell$ である。
巡回代数のノルム判定\cite[\S1, equation~(2)]{MinacSchultzSwallowDescent39}により、
$[(\Xi,U)_t]^a=0$ なら $U^a=N_{E/K_\nu}(b)$ でなければならない。
$E$ の付値を $\nu_E(x)=1$ と規格化すると、
完全分岐性より $\nu(N(b))=\nu_E(b)$ である。
従って $b=x^a b_0$、$b_0$ は付値環の単元と書ける。
\[
 N(x)=(-1)^{\ell-1}\Xi
\]
なので、ノルム等式を単元部分で剰余体に移すと
\[
 r^a\in\kappa_\nu^{\times\ell}
\]
が必要となる。符号は $\mathbb C$ 内で $\ell$ 乗根を持つため無関係である。
完全分岐拡大では全ての共役の剰余が一致するので、単元のノルムの剰余は
確かに $\ell$ 乗である。
しかし $r$ に沿う付値を使えば、$r^a$ が $\ell$ 乗であるためには
$\ell\mid a$ が必要である。
従って巡回代数のBrauer周期は、$K_\nu$ 上でも $K_0$ 上でも正確に $\ell$ である。
指数は周期の倍数かつ代数の次数 $\ell$ の約数なので、指数も $\ell$ となる。
これは体拡大後の代数が依然として斜体であることも証明する。

最後に、反射的加群 $M$ の一般点は非零有限次元 $K_0$ ベクトル空間である。
その自己準同型環は分裂行列環になるので、上で得た非零Brauer類と両立しない。
中心線形Morita同値も一般点のBrauer類を保存するため、同じ結論となる。
\end{proof}

この完備局所障害は、$\Xi\ne0$ にある閉点の完備化と矛盾しない。
後者では $\Xi$ は非零定数項を持つ単元なので、Henselの補題で
$\ell$ 乗根を取り、巡回代数を分裂させられる。
最深点では $\Xi,U$ がともに消え、式\eqref{eq:deepest-valuation-weights-v40}の
付値がその障害を検出する。
本定理は、中心自身が別のNCCRを持たないという主張ではない。
任意の二変数混合項に対する全中心の決定は、依然として別の問題である。

\begin{corollary}[$Q$ で割り切れる混合項の最深点]
\label{cor:deepest-Q-divisible-v40}
$f\in Q\mathbb C[Q,u]$ とし、
$Z_f=Z(S_f)$ の原点 $o_f=V(Q,J,\Xi,\Omega,U,V)$ を取る。
この点での中心ファイバーは $\mathcal F_\ell$ と同型である。
また、$\widehat Z_f$ は整域であり、その分数体上の $\widehat S_f$ の
Brauer周期と指数はともに $\ell$ である。
従って、この完備代数にも定理\ref{thm:deepest-Brauer-v40}と同じ
自己準同型環および中心線形Morita同値についての否定的結論が成立する。
\end{corollary}

\begin{proof}
$Q$ は中心元なので、$f=Qg$ と書けば $\delta_f=Q\delta_g$ である。
従って $Q=0$ を課すとOre関係は $S_0$ のものになり、
$J=xw+f$ も $xw$ になる。他の四つの中心座標は生成元の $\ell$ 乗で変わらない。
これで最深ファイバーが同じ $\mathcal F_\ell$ となる。

完備化について、$f$ の $Q$ の一次項を $cQ$ と書き、$J'=J-cQ$ と置く。
これは $cQ$ が中心的であるため、同じ代数の中心座標の変更である。
$f-cQ$ に現れる単項式 $Q^a u^k$ は全て
\[
 a\ge1,\qquad a+\frac{k}{\ell}>1
\]
を満たす。$\operatorname{wt}(u)=1/\ell$、
$\operatorname{wt}(Q)=\operatorname{wt}(J')=1$ として巡回積を展開すれば、
\[
 N_{f-cQ}(Q,U,J')=(J')^\ell+\text{重みが $\ell$ より大きい項}
\]
となる。巡回不変性によって $u$ の指数は $\ell$ の倍数なので、
この重みは中心上の整数重み $\operatorname{wt}(U)=1$ と一致する。
式\eqref{eq:deepest-valuation-weights-v40}の重みで、二つの中心方程式の初期形は再び
\[
 \Xi\Omega-(J')^\ell,\qquad UV-c_\ell\Xi^2,
 \qquad c_\ell=(-1)^\ell\beta^\ell
\]
である。これらは多項式環内で正則列を成す。
実際、第一式による商は整域であり、第二式はその非零元である。
従って、完備環の商を一式ずつ取る際に初期形が非零因子であることを用いると、
付随次数環は正確にこの二式による商となる。
これは式\eqref{eq:deepest-associated-graded-v40}と同じ整域である。
よって同じ離散付値、剰余体 $\mathbb C(q,j,r)$、Eisenstein拡大が得られる。
一般点の巡回表示も $(\Xi,U)_t$ のままなので、
定理\ref{thm:deepest-Brauer-v40}のノルム判定をそのまま適用できる。
\end{proof}

この系は最深点についての主張であり、一般の $f\in Q\mathbb C[Q,u]$ の
全ファイバーが $S_0$ と同じだと主張するものではない。
次の例では、実際に全中心の特異部分が変わる。

\subsection{同じ関数斜体を持つが全中心が異なる二つの量子化}

\begin{proposition}[標識を忘れた全中心の非同型性]
\label{prop:unmarked-global-centers-v40}
$m=4$、$\ell=5$、$\beta=1$ とし、$f_0=0$、$f_1=Qu^3$ を取る。
二つの中心 $Z_i=Z(S_{f_i})$ は
\begin{align*}
 Z_0&=\mathbb C[Q,J,\Xi,\Omega,U,V]/
       (\Xi\Omega-J^5,\ UV+Q^5+\Xi^2),\\
 Z_1&=\mathbb C[Q,J,\Xi,\Omega,U,V]/
       (\Xi\Omega-J^5+Q^5U^3,\ UV+Q^5+\Xi^2)
\end{align*}
と表示される。$Z_0$ の特異部分は次元二である一方、
$Z_1$ の被約特異部分は
\begin{equation}\label{eq:unmarked-singular-axes-v40}
 Q=J=\Xi=0,\qquad UV=U\Omega=V\Omega=0
\end{equation}
という三本の座標軸の合併であり、次元一である。
従って、標識を忘れても $Z_0\not\simeq Z_1$ である。
代数 $S_{f_0},S_{f_1}$ は同じ関数斜体と同型な $D(\Xi)$ 上の局所化を持つが、
次数を忘れた $\mathbb C$ 代数としてMorita同値ではなく、dg導来Morita同値でもない。
両者の最深ファイバーは同じ $\mathcal F_5$ で、その完備局所Brauer周期・指数もともに $5$ である。
\end{proposition}

\begin{proof}
$Qu^3$ は元の重みで次数五を持ち、$m=4$ に対応する混合項である。
巡回軌道の長さは五なので $N_{Qu^3}=J^5-Q^5U^3$ となり、中心の表示を得る。
$Z_0$ の特異部分の次元は定理\ref{thm:global-fibers-v40}から二である。
$Z_1$ の二つの方程式のJacobianは、列を $(Q,J,\Xi,\Omega,U,V)$ の順に取ると
\[
 \begin{pmatrix}
 5Q^4U^3&-5J^4&\Omega&\Xi&3Q^5U^2&0\\
 5Q^4&0&2\Xi&0&V&U
 \end{pmatrix}
\]
である。$\Xi\ne0$ では $\Omega,\Xi$ の二列が独立である。
$\Xi=0,J\ne0$ では第一の方程式から $Q,U\ne0$ となり、$J,V$ の二列が独立である。
残る $\Xi=J=0$ では二つの方程式から $Q=0,UV=0$ を得る。
このとき二行はそれぞれ $\Xi$ 列の $\Omega$ と、$U,V$ 列の $(V,U)$ だけを持つので、
階数が二未満となる条件は正確に $\Omega=0$ または $U=V=0$ である。
これが式\eqref{eq:unmarked-singular-axes-v40}である。

どちらの代数も $\Xi$ を可逆にすれば
$w=x^{-1}(J-f_i)$ によって $B[x^{-1}][J]$ と同型になる。
従ってこの局所化と関数斜体は同じである。
一方、Morita同値は中心を保存する。
さらにdg導来Morita同値はHochschildコホモロジー代数を保存し
\cite[Theorem~3.2 and its corollary]{KellerDerived2003}、
その次数零部分は通常の代数の中心である。
中心の非同型性から、いずれの同値も存在しない。
最深点の二つの主張は系\ref{cor:deepest-Q-divisible-v40}による。
\end{proof}

これはアフィン代数とそのdg導来Morita型についての区別である。
次数付き商圏 $\operatorname{qgr}S_{f_i}$ やその射影的導来圏の非同値までは結論していない。

\section{境界の形式McKay模型と例外因子上のBrauer留数}
\label{sec:boundary-ramification-v40}

前節の巡回被覆は、単純加群を区別するための被覆であった。
本節では同じ被覆が、中心の可換解消上で量子化をAzumaya代数へ延長する障害も測ることを示す。
閉点の完成と、境界の一般点での完成を区別することが必要である。
前者では巡回被覆が分裂するが、後者ではその関数体が残る。

\subsection{正規化境界に沿う全形式代数}

第\ref{sec:endpoint-boundary-v39}節の $C_H=S_H[U^{-1}]$ を用いる。
$h$ の安定化群の位数を $g$、$n=\ell/g$、$z=u^g$ と置く。
被約スペクトル因子の関数体は $F=\mathbb C(Q,z)$ であり、$U=z^n$ である。
巡回作用により得られる相異なる $n$ 個の元を
\[
 h_i=h(Q,t^i u)\in F\quad(0\le i<n)
\]
とする。ここで $h_0=h$ とし、$h_n=h_0$ と読む。
以下で用いる積は
\[
 s=J-h_0,\qquad b_i=s+h_0-h_i,\qquad
 E(s)=\prod_{i=1}^{n-1}b_i.
\]
$E(0)\ne0$ であり、$n=1$ の場合は $E=1$ とする。
第\ref{sec:endpoint-boundary-v39}節のノルムは $(sE(s))^g$ となる。
従って $\xi=\Xi$、$\upsilon=\Omega/E(s)^g$ によって、
一般境界点の完成中心は
\begin{equation}\label{eq:boundary-transverse-center-v40}
 R_F=\frac{F[[\xi,\upsilon,s]]}{(\xi\upsilon-s^g)}
\end{equation}
となる。これは元の中心の余次元二の点における完成である。
係数体の埋込みは次のように選ぶ。剰余体における $z$ を、
$z^n=U$ の根としてHenselの補題で持ち上げる。
$nz^{n-1}$ は単元なのでこの持上げは一意であり、元の $Q$ とともに
$F=\mathbb C(Q,z)$ の埋込みを定める。

$L=F[\kappa]/(\kappa^g-z)$ と置き、
$\sigma(\kappa)=t^{-n}\kappa$ とする。$L/F$ は巡回次数 $g$ の体拡大である。
次の代数をその中心イデアル $(P^g,R^g,PR)$ で完備化する：
\begin{equation}\label{eq:boundary-formal-corner-v40}
 \mathcal T_g=
 \widehat{\frac{L\langle P,R\rangle}
 {(Pa-\sigma(a)P,\ Ra-\sigma^{-1}(a)R,\ PR-RP\ :\ a\in L)}}.
\end{equation}
その中心からの作用を $\xi=P^g$、$\upsilon=R^g$、$s=PR$ で定める。
$g=1$ では $\mathcal T_1=F[[P,R]]$ と解釈する。

\begin{theorem}[巡回降下の全形式延長]
\label{thm:boundary-formal-descent-v40}
一般境界点における完成代数は
\begin{equation}\label{eq:boundary-formal-matrix-v40}
 C_H\otimes_{Z(C_H)}R_F\simeq M_n(\mathcal T_g)
\end{equation}
である。$\mathcal T_g$ の中心は正確に $R_F$ であり、中心上の階数は $g^2$ である。
さらに、係数を $L$ へ拡大すると
\[
 \mathcal T_g\widehat\otimes_F L
 \simeq L[[a,b]]\#\mu_g,
\]
ここで $\zeta=t^n$ として、群の生成元は $(a,b)\mapsto(\zeta a,\zeta^{-1}b)$ と作用する。
\end{theorem}

\begin{proof}
$u^g/z$ の相異なる $n$ 個の固有値により冪等元が得られる。
$b_1,\ldots,b_{n-1}$ は完成環上で可逆なので、対応する往復矢を縮約できる。
縮約後の一つの角で $\kappa=u$、$P=x^n$、$R=E(s)^{-1}W^n$ と置く。
一般化Weyl関係により
\[
 \kappa^g=z,\quad P\kappa=t^{-n}\kappa P,\quad
 R\kappa=t^n\kappa R,\quad PR=RP=s,
\]
および $P^g=\xi$、$R^g=\upsilon$ を得る。

逆写像を明示できる。$n\times n$ 行列単位を $E_{ij}$ とし、
\begin{align*}
 u&=\operatorname{diag}(\kappa,t\kappa,\ldots,t^{n-1}\kappa),\\
 x&=\sum_{i=0}^{n-2}E_{i+1,i}+P E_{0,n-1},\\
 W&=\sum_{i=1}^{n-1}b_iE_{i-1,i}+R E_{n-1,0}
\end{align*}
とする。$n=1$ でも和を空として同じ式を用いる。
これらは $xu=t^{-1}ux$、$Wu=tuW$、
$xW=\operatorname{diag}(b_0,\ldots,b_{n-1})$ と
$Wx=\operatorname{diag}(b_1,\ldots,b_{n-1},b_0)$ を満たす。
また $x^n=P I$、$W^n=E(s)R I$、$u^\ell=z^n I$ である。
固有冪等元と可逆な内部矢から全ての行列単位が得られ、角には $\kappa,P,R$ が残る。
従って写像は全射である。

$L$ 上の単項式 $P^aR^b$ による正規形から、
中心は $F$ 係数で $a-b\equiv0\pmod g$ の単項式で生成される。
よって中心は $F[[P^g,R^g,PR]]=R_F$ である。
一般点では $P$ が可逆で $R=sP^{-1}$ と書け、$\kappa^iP^j$、$0\le i,j<g$ が基底となる。
従って上の全射は一般点で同型である。
元の一般化Weyl代数は中心上捩れなしであり、局所化と完成の平坦性により
左辺も $R_F$ 上捩れなしである。従って核は零である。

最後に係数拡大した $L$ における $z$ の根を $\kappa_0$ と書く。
元 $c=\kappa/\kappa_0$ は $c^g=1$、$cP=\zeta Pc$、$cR=\zeta^{-1}Rc$ を満たす。
正規形を比較すれば、$a=P$、$b=R$、$c\in\mu_g$ とするskew group表示を得る。
\end{proof}

\begin{corollary}[一般境界の閉点における非可換クレパント解消]
\label{cor:boundary-closed-McKay-v40}
追加の値の一致を除いた正規化境界の複素閉点 $p$ では、
適当な形式変数により
\[
 \widehat{C}_{H,p}\simeq
 M_n\bigl(\mathbb C[[a,b,q_1,q_2]]\#\mu_g\bigr),
\]
ここで $\mu_g$ は $q_1,q_2$ を固定する。
これは完成中心の非可換クレパント解消である。
特に横断特異点は $A_{g-1}$ 型であり、$g=1$ では滑らかである。
\end{corollary}

\begin{proof}
閉点では $z\ne0$ なので、Henselの補題により完成係数環で $z$ の $g$ 乗根を選べる。
前定理の構成を $\mathbb C[[q_1,q_2]]$ 上で行い、$s=ab$ を消去すれば表示を得る。
群作用は行列式一で、非自明な固定部分の余次元は二である。
Auslanderの同型とskew group代数の標準的な議論により
\[
 \mathbb C[[a,b,q_1,q_2]]\#\mu_g
 \simeq\operatorname{End}_{\mathbb C[[a,b,q_1,q_2]]^{\mu_g}}
 \bigl(\mathbb C[[a,b,q_1,q_2]]\bigr)
\]
は最大Cohen--Macaulayで大域次元四である。
この商特異点に対する非可換クレパント解消は
Van den Bergh\cite[Example~1.1]{VanDenBerghMcKay38}の標準例の完成版である。
行列環を取っても同じ性質を保つ。
\end{proof}

\subsection{可換解消へ移ったときの分岐}

$g>1$ とする。$K_F=\operatorname{Frac}(R_F)$ 上の角代数は
\[
 D_g=(L\otimes_F K_F/K_F,\sigma,\xi)
\]
という巡回代数である。
ここで巡回代数の規約は、生成元 $P$ が $Pa=\sigma(a)P$、$P^g=\xi$ を満たすものとする。
$\chi\in H^1(F,\mathbb Z/g)$ は $L/F$ に対応し、$\chi(\sigma^{-1})=1$ となる指標とする。
この逆元の指定は符号を固定するためである。
実際 $aP=P\sigma^{-1}(a)$ なので、
Gille--Szamuely\cite[Propositions~2.5.2 and~4.7.3]{GilleSzamuelyBrauer40}の規約で
そのBrauer類は $\chi\cup(\xi)$ である。

$\widetilde X_F\to\operatorname{Spec}R_F$ を横断 $A_{g-1}$ 特異点の最小解消とする。
例外曲線を、対応する付値が
\begin{equation}\label{eq:boundary-exceptional-valuations-v40}
 v_i(\xi)=i,\qquad v_i(\upsilon)=g-i,\qquad v_i(s)=1
 \quad(1\le i<g)
\end{equation}
となるように $E_i$ と番号付ける。
これらの式は、解消の $k$ 番目のトーリックchart
\[
 \xi=a_k^{k+1}b_k^k,\qquad
 \upsilon=a_k^{g-k-1}b_k^{g-k},\qquad s=a_kb_k
 \quad(0\le k<g)
\]
からも直接得られる。$\kappa(E_i)=F(t_i)$ である。

\begin{theorem}[単純加群の被覆とBrauer留数]
\label{thm:boundary-residues-v40}
各例外因子におけるBrauer留数は
\begin{equation}\label{eq:boundary-residues-v40}
 \partial_{E_i}[D_g]=i\,\chi|_{F(t_i)},\qquad
 \operatorname{ord}\bigl(\partial_{E_i}[D_g]\bigr)=\frac{g}{\gcd(g,i)}.
\end{equation}
従って $D_g$ の周期と指数はともに $g$ である。
このBrauer類は最小可換解消上のAzumaya代数の類へは延長できず、
その解消を支配する正則な解消でも同様である。

一方、係数拡大 $L/F$ により全ての留数は消える。
有限係数拡大 $F'/F$ が全ての例外留数を消すなら、$F'$ は $z$ の $g$ 乗根を含み、
$[F':F]$ は $g$ の倍数である。
従って必要な最小係数被覆は、境界上の単純加群を分離する巡回被覆そのものである。
\end{theorem}

\begin{proof}
指標 $\chi$ は各 $v_i$ に関して不分岐である。
巡回代数の留数公式
$\partial_v(\chi\cup b)=v(b)\overline\chi$
\cite[Proposition~5.8]{HarariBrauer40}と
式\eqref{eq:boundary-exceptional-valuations-v40}から最初の等式を得る。
$z$ は $F=\mathbb C(Q,z)$ の独立変数であり、$z$ に関する付値により
$\kappa^g-z$ はEisenstein多項式である。
さらに純超越拡大 $F(t_i)/F$ はこの巡回拡大の次数を変えない。
よって $\chi|_{F(t_i)}$ の位数は $g$ であり、留数の位数も従う。

$E_1$ の留数が位数 $g$ なので $[D_g]$ の周期は少なくとも $g$ である。
$D_g$ は次数 $g$ の巡回代数であるから、周期は指数を割り、指数は $g$ を割る。
従って両者は $g$ に等しい。
Azumaya代数の一般点の類は各余次元一の点で不分岐である
\cite[Proposition~4.23]{HarariBrauer40}。
非零の留数が延長を妨げる。支配する解消においても $E_1$ の狭義変換の付値は変わらない。

係数拡大で $\chi$ を零にすれば巡回代数自体が分裂するので、十分性は明らかである。
逆に $E_1$ の留数が零なら、$\chi$ の $F'(t_1)$ への制限は零である。
有限定数拡大の指標は純超越拡大への制限で単射に写るので、$F'$ への制限も零である。
これは $z$ が $F'$ で $g$ 乗となることと同値であり、$L$ が $F'$ に埋め込まれる。
\end{proof}

例えば $g=6$ なら、五つの例外曲線における留数の位数は
\[
 (6,3,2,3,6)
\]
である。横断中心の式だけは係数体上ですでに $A_5$ 型であるが、
量子化の類にはこの分岐データが残る。
また $\mathbb C(Q,U)\subset\mathbb C(Q,z)\subset\mathbb C(Q,\kappa)$ の次数は
$n,g$ で、積は $\ell$ である。
スペクトル因子を正規化する次数と、その上の単純加群を分離する次数が、
前節の一般点の最小分裂次数 $\ell$ をこの形で分解する。

以上は $U\ne0$ の一般境界、すなわち追加衝突を除いた部分の主張である。
任意の衝突点の形式型や、全中心の可換解消を一括して分類するものではない。
また閉点完成の非可換クレパント解消と、一般境界点での非零Brauer留数は両立する。
前者では係数 $z$ の根が存在し、後者ではその根が巡回拡大 $L/F$ を与えるためである。

\section{証明の範囲と残る境界の問題}

正共鳴の全整数にわたる結果は、指数半群、Ore表示、正規スライス、有限整拡大による。
原点根基の冪零指数には全重みでの所属恒等式と非所属評価を与えた。
量子・Poisson素スペクトルの一致は、正規元の内部自己同型から局所冪零導分を回収する補間による。
葉の閉包公式は全域で成り立つが、無限次元単純加群そのものの全同型類や部分解消上の代数層は別の問題である。

一の冪根では、巡回ノルム、有限エタール被覆上の重み冪等元、非零辺の縮約により全位数を扱った。
行列代数になる開部分では、閉点ファイバーの単純性に加えて局所自由性も示す。
非Azumayaファイバーの全体は、道の簡約による上界と明示的な行列表現による下界を合わせて決定する。
特異部分への二次降下と一般の巡回降下も、互いに逆な行列表現を構成して証明した。
計算コードは有限個の位数についてこれらの全関係を厳密に照合するが、
全位数への結論には本文の一般証明を用いる。

Brauer周期と指数は、整域性と巡回ノルムの付値によって決まり、有限位数の標本からの推測ではない。
分裂被覆上で得られるのはベクトル束の自己準同型環である。
元の中心上では非零のBrauer類が残るため、通常の自己準同型環としてのNCCRと
同一視できない。この障害は、既知の中心開部分の閉点での完成局所McKay模型と両立する。
さらに全中心の最深点では、完成後にも指数 $\ell$ が残ることを独立した付値計算で証明した。

$U\ne0$ の中心は混合項 $H$ の剰余 $h$ をノルムとして記憶する。
この復元の結論は $Q,U,J$ を固定したものである。
中心元を $H$ に加える操作はOre導分を変えず、$J$ の移動に対応する。
全中心を求めた $f=0$ と $f=Qu^3$ の例では、特異部分の次元から
標識を忘れた全中心の非同型性も従う。
ここからアフィン代数のMorita・dg導来Morita非同値を得るが、射影導来圏の非同値までは結論しない。

混合項が $f(Q,u)+xG$ となる族については全中心を決定した。
混合項が中心元と $xG$ の和になる部分族では全閉点ファイバーも分類し、
$f\in Q\mathbb C[Q,u]$ では最深点のファイバーと完備局所Brauer障害が同じであることを示した。
一般境界の全形式模型からは、単純加群の巡回被覆と例外因子のBrauer留数を対応させた。
ここで用いるMcKay/NCCRの一般理論と巡回代数の留数公式は既存の理論であり、
本稿の内容はこの量子化族における明示表示と、その大域・完備局所的な相違である。

残る対象は、この一変数形に帰着しない混合項 $H$ に対する全中心と深い境界、
追加衝突点の全形式代数、および標識を忘れた射影圏の同値である。
最深点の代数自身に対するEnd型の障害は、中心が別のNCCRを持たないという結論ではない。
中心局所化や係数体上のBrauer障害を、任意の導来埋込みの不存在へ読み替えない。

\begin{thebibliography}{99}
\interlinepenalty=10000
\bibitem{VKPoisson34} 本研究、公表原稿IX、
\emph{重み付きPoisson変形の非被約成分とCalabi--Yau量子化}, v34.
\nolinkurl{weighted_poisson_components_ja_v34.pdf}.
\bibitem{VKCategories36} 本研究、公表原稿X、
\emph{重み付き量子射影圏の変形環と対称性}, v40.
\nolinkurl{quantum_categories_ja_v40.pdf}.
\bibitem{LecoutreSierra2019}
C.~Lecoutre and S.~J.~Sierra,
\emph{A new family of Poisson algebras and their deformations},
Nagoya Math. J. \textbf{233} (2019), 32--86,
\href{https://doi.org/10.1017/nmj.2017.29}{doi:10.1017/nmj.2017.29};
\href{https://arxiv.org/abs/1611.04306}{arXiv:1611.04306}.
\bibitem{ORS2013}
J.~\"Oinert, J.~Richter and S.~Silvestrov,
\emph{Maximal commutative subrings and simplicity of Ore extensions},
Journal of Algebra and Its Applications \textbf{12} (2013), no.~4, 1250192, 16 pp.
\url{https://doi.org/10.1142/S0219498812501927}.
\url{https://arxiv.org/abs/1111.1292}.

\bibitem{JordanOre2014}
D.~A.~Jordan, \emph{Ore extensions and Poisson algebras},
Glasgow Mathematical Journal \textbf{56} (2014), no.~2, 355--368.
\url{https://doi.org/10.1017/S0017089513000293}.
\url{https://arxiv.org/abs/1212.4063}.

\bibitem{BellWangYee2019}
J.~P.~Bell, X.~Wang, and D.~Yee,
\emph{The Dixmier--Moeglin equivalence, Morita equivalence, and homeomorphism of spectra},
arXiv:1807.11813 (2018), proof of Theorem~1.1.
\url{https://arxiv.org/abs/1807.11813}.

\bibitem{SuarezVivas2015Root38}
M.~Su\'arez-Alvarez and Q.~Vivas,
\emph{Automorphisms and isomorphisms of quantum generalized Weyl algebras},
J. Algebra \textbf{424} (2015), 540--552,
\href{https://doi.org/10.1016/j.jalgebra.2014.08.045}{doi:10.1016/j.jalgebra.2014.08.045},
\href{https://arxiv.org/abs/1206.4417}{arXiv:1206.4417}.

\bibitem{VanDenBerghMcKay38}
M.~Van den Bergh,
\emph{Non-commutative crepant resolutions},
in \emph{The Legacy of Niels Henrik Abel}, Springer, 2004, 749--770;
corrected version,
arXiv:\href{https://arxiv.org/abs/math/0211064}{math/0211064}, 2009.

\bibitem{MinacSchultzSwallowDescent39}
J.~Min\'{a}\v{c}, A.~Schultz and J.~Swallow,
\emph{Cyclic algebras, Schur indices, norms, and Galois modules},
Annales des sciences math\'{e}matiques du Qu\'{e}bec \textbf{35} (2011), no.~1, 123--136.
\href{https://arxiv.org/abs/math/0410536}{arXiv:math/0410536};
the splitting--norm criterion is recalled in \S1, equation~(2).

\bibitem{DGOBoundary39}
Yu.~A.~Drozd, B.~L.~Guzner and S.~A.~Ovsienko,
\emph{Weight modules over generalized Weyl algebras},
Journal of Algebra \textbf{184} (1996), 491--504.
\url{https://www.imath.kiev.ua/~drozd/Guzner.pdf}.
\bibitem{CGJBoundary39}
N.~Cooney, I.~Ganev and D.~Jordan,
\emph{Quantum Weyl algebras and reflection equation algebras at a root of unity},
Journal of Pure and Applied Algebra \textbf{224} (2020), no.~12, 106440.
\href{https://arxiv.org/abs/1907.11005}{arXiv:1907.11005};
doi:10.1016/j.jpaa.2020.106440.

\bibitem{HarariBrauer40}
D.~Harari, \emph{The Brauer group of fields and schemes}, M2 lecture notes, 2023--2024.
Propositions~4.23 and~5.8.
\url{https://www.imo.universite-paris-saclay.fr/~david.harari/enseignement/m2brauer24/brauer24.pdf}.

\bibitem{GilleSzamuelyBrauer40}
P.~Gille and T.~Szamuely, \emph{Central Simple Algebras and Galois Cohomology},
Cambridge Studies in Advanced Mathematics 101, Cambridge University Press, 2006.
Propositions~2.5.2 and~4.7.3.
\url{https://www.math.ens.psl.eu/~benoist/refs/Gille-Szamuely.pdf}.

\bibitem{KellerDerived2003}
B.~Keller, \emph{Derived invariance of higher structures on the Hochschild complex},
preprint, August 2003, revised May 16, 2018, Theorem~3.2 and its corollary.
\url{https://webusers.imj-prg.fr/~bernhard.keller/publ/dih.pdf}.

\end{thebibliography}

\end{document}
